% vim: foldmethod=marker foldmarker=<<<,>>>
\ifx\danslelivre\undefined
\documentclass[10pt]{../fichiers-configuration-latex/smfart}
\usepackage{ifthen}
\newboolean{english}
\newboolean{francais}
\setboolean{english}{false}
\setboolean{francais}{true}
\input{../fichiers-configuration-latex/paquets}
\selectlanguage{francais}
\input{../fichiers-configuration-latex/lemme-proposition-theoreme}
\input{../fichiers-configuration-latex/numerotation}
\input{../fichiers-configuration-latex/macros}

\author{Joël Riou}
\address{Université Paris-Sud 11, Bât.~425, 91405 Orsay, France}
\email{joel.riou@math.u-psud.fr}
\title{XVI. Classes de Chern, morphismes de Gysin, pureté absolue}

\date{exposés oraux des 16 juin 2006 et 13 avril 2007, notes du \today}

\setcounter{tocdepth}{3}

\begin{document}

\begin{center}
XVI. Classes de Chern, morphismes de Gysin, pureté absolue\\
Joël Riou\\
exposés oraux des 16 juin 2006 et 13 avril 2007\\
\end{center}

\begin{center}
\verb$Id: gysin.tex,v 1.98 2009-08-20 10:06:45 cvs Exp $
\end{center}

\tableofcontents

\else
\setboolean{english}{false}
\setboolean{francais}{true}
\selectlanguage{francais}
\chapter{Classes de Chern, morphismes de Gysin, pureté absolue}
\addtocontents{toc}{par Joël Riou}
\begin{center}
Joël Riou
\end{center}
\fi

Dans ces notes, on présente la nouvelle démonstration par Ofer Gabber du
théorème de pureté cohomologique absolue, annoncée dans
\cite{ABELIEN@Gabber}. La
section~\ref{16-section-classes-chern} rappelle la construction des classes de
Chern en cohomologie étale. Celles-ci servent dans la
section~\ref{16-section-gysin} qui consiste en la construction et l'étude des
propriétés des morphismes de Gysin associés aux morphismes d'intersection
complète lissifiables. Dans la section~\ref{16-section-purete}, ces morphismes
de Gysin sont utilisés pour donner une formulation précise du théorème de
pureté absolue (théorème~\ref{16-theoreme-purete-absolue}). La démonstration
du théorème de pureté (différente de celle rédigée dans
\cite{AZUMINO@Fujiwara})
s'appuie notamment sur les résultats de géométrie logarithmique établis
dans l'exposé VI.

\bigskip

Dans tout cet exposé, on fixe un entier naturel $n\geq 1$. Tous les schémas
seront supposés être des schémas sur $\SP\left(\ZZ\cro{\frac 1 n}\right)$. On note
$\Lambda$ le faisceau d'anneaux constant de valeur $\ZZ/n\ZZ$, $\Lambda(1)$
le faisceau des racines $n$-ièmes de l'unité (pour la topologie étale)
et $\Lambda(r)$ ses
puissances tensorielles, auxquelles on peut donner un sens pour tout
$r\in\ZZ$.

\section{Classes de Chern}%<<<1
\label{16-section-classes-chern}

Dans cette section, on rappelle la construction des classes de Chern
des fibrés vectoriels sur des schémas généraux à valeurs dans la 
cohomologie étale.
On s'appuie sur le calcul de la cohomologie étale des fibrés projectifs de
\sga{5}{vii}{2} et sur la méthode de \cite{CHERN@Grothendieck}.
Les démonstrations sont parfois différentes de celles
de \sga{5}{vii}{3} : on s'est efforcé de donner une présentation aussi
\guil{économique} que possible.

À la différence de l'exposé oral qui utilisait un langage géométrique,
dans ces notes, un fibré vectoriel est un Module $\mathscr E$
localement libre de rang fini et le fibré projectif de $\mathscr E$ est
le fibré des hyperplans défini dans \ega{ii}{4.1.1} :
$\PP(\mathscr E)=\Proj(\SS^\star\mathscr E)$ où $\SS^\star\mathscr E$
est l'Algèbre symétrique de $\mathscr E$.

\begin{definition}\label{16-definition-c-1}
Soit $X$ un $\ZZn$-schéma. Soit $\mathscr L$ un fibré en droites sur $X$.
Le faisceau des sections inversibles de $\mathscr L$ est naturellement
muni d'une structure de torseur sous le schéma en groupes $\Gm$. La classe
d'isomorphisme de $\mathscr L$ définit donc un élément dans $H^1_\et(X,\Gm)$.
On note $c_1(L)\in \HH^2_\et(X,\Lambda(1))$ l'image de cet élément par le
morphisme de bord déduit de la suite exacte courte de Kummer :
\[0\to \Lambda(1)\to \Gm\vers {\cro n}\Gm \to 0\;\footnote{%
Une grande liberté d'interprétation est laissée à l'imagination du lecteur
au sujet des conventions de signes à utiliser dans cette définition de
$c_1$. La même liberté me semblant exister dans \sga{4}{xviii}{1.1}
pour le morphisme
trace, il me paraît illusoire de faire davantage de zèle. Cependant, à un
endroit de ces notes, ce choix ne sera pas neutre : il faut que les choix
de signes soient mutuellement cohérents de sorte que l'on ait une
compatibilité entre le degré des fibrés en droites sur la droite
projective, le morphisme trace et la première classe de Chern (cf.
démonstration du lemme~\ref{16-lemme-compatibilite-trace-c-un}).}\text{.}\]
\end{definition}

Si $\mathscr L$ et $\mathscr L'$ sont deux 
fibrés en droites sur $X$, on a la relation
d'additivité :
\[c_1(\mathscr L\otimes \mathscr L')=
c_1(\mathscr L)+c_1(\mathscr L')\in \HH^2(X,\Lambda(1))\;%
\footnote{Il existe des théories cohomologiques
\guil{orientées} pour lesquelles cette propriété de la première classe de
Chern n'est pas satisfaite, cf.~\cite{CA1@Morel-Levine}.}\text{.}\]

Notons que les classes de Chern de fibrés en droites résident dans les degrés
pairs de la cohomologie étale, elles commutent donc avec toutes les classes
de cohomologie. Notons
aussi que si $f\colon Y\to X$ est un morphisme et 
$\mathscr L$ un fibré en droites
sur $X$, alors $f^\star(c_1(\mathscr L))=c_1(f^\star \mathscr L)$.

\begin{theoreme}[Formule du fibré projectif]%
\label{16-theoreme-formule-du-fibre-projectif}
Soit $X$ un $\ZZn$-schéma. Soit $\mathscr E$ un fibré 
vectoriel de rang constant $r$
sur $X$.
On note $\pi\colon\PP(\mathscr E)\to X$
le fibré projectif de $\mathscr E$.
On pose $\xi=c_1(\OO(1))\in \HH^2(X,\Lambda(1))$\;\footnote{Le faisceau 
$\OO(1)$ est le faisceau fondamental sur $\PP(\mathscr E)$ : c'est le 
quotient inversible de $\pi^\star \mathscr E$ par 
l'hyperplan universel.}.
Alors, les puissances $\xi^i\in \HH^{2i}(X,\Lambda(i))$ de $\xi$
définissent un isomorphisme dans $\detaleplus X$ :
\[(1,\xi,\dots,\xi^{r-1})\colon
\bigoplus_{i=0}^{r-1} \Lambda(-i)[-2i]
\iso \R\pi_\star \Lambda\]
\end{theoreme}

D'après le théorème de changement de base pour un morphisme propre, on peut
supposer que $X$ est le spectre d'un corps algébriquement clos $k$. On se
ramène ainsi au calcul de l'algèbre de cohomologie étale des espaces
projectifs sur $k$, cf.~\sga{5}{vii}{2}.

\begin{theoreme}\label{16-theoreme-classes-de-chern}
Il existe une unique manière de définir, pour tout $\ZZn$-schéma $X$ et
tout fibré vectoriel $\mathscr E$, 
des éléments $c_i(\mathscr E)\in \HH^{2i}_\et(X,\Lambda(i))$
pour tout $i\in\NN$ de sorte que si l'on définit la série formelle
$c_t(\mathscr E)=\sum_{i\geq 0} c_i(\mathscr E)t^i$,
on ait les propriétés suivantes :
\begin{itemize}
\item la série formelle $c_t(\mathscr E)$ ne dépend que 
de la classe d'isomorphisme
du fibré vectoriel $\mathscr E$ sur le $\ZZn$-schéma $X$ ;
\item si $f\colon Y\to X$ est un morphisme de $\ZZn$-schémas
et $\mathscr E$ un fibré vectoriel sur
$X$, alors $f^\star(c_t(\mathscr E))=c_t(f^\star \mathscr E)$ ;
\item si $0\to \mathscr E'\to \mathscr E\to \mathscr E''\to 0$ est une
suite exacte courte de fibrés vectoriels sur un $\ZZn$-schéma $X$, on a la
relation de Cartan-Whitney :
\[c_t(\mathscr E)=c_t(\mathscr E')c_t(\mathscr E'')\;\text{;}\]
\item si $\mathscr L$ est un fibré en droites 
sur un $\ZZn$-schéma $X$, la classe
$c_1(\mathscr L)$ est celle de la définition~\ref{16-definition-c-1} et 
\[c_t(\mathscr L)=1+c_1(\mathscr L)t\;\text{.}\]
\end{itemize}

On a alors les relations $c_0(\mathscr E)=1$ et 
$c_i(\mathscr E)=0$ pour $i>\rang \mathscr E$ pour tout
fibré vectoriel $\mathscr E$ sur un $\ZZn$-schéma $X$.
\end{theoreme}

La démonstration utilise plusieurs constructions géométriques :

\begin{proposition}[Principe de scindage I]\label{16-proposition-scindage-1}
Soit $X$ un $\ZZn$-schéma. Soit $\mathscr E$ un 
fibré vectoriel de rang $r$. On note
$\pi\colon \Drap(\mathscr E)\to X$ le fibré des drapeaux complets de
$\mathscr E$. Les
propriétés suivantes sont satisfaites :
\begin{itemize}
\item le fibré vectoriel $\pi^\star \mathscr E$ admet une filtration (canonique)
$0=\mathscr M_0\subset \mathscr M_1\subset \dots\subset \mathscr M_r=
\pi^\star \mathscr E$ par des fibrés
vectoriels de sorte que pour tout entier $1\leq i\leq r$, le quotient
$\mathscr L_i=\mathscr M_i/\mathscr M_{i-1}$ soit un fibré en droites ;
\item le morphisme canonique $\Lambda\to \R\pi_\star\Lambda$ est un
monomorphisme scindé dans $\detaleplus X$.
\end{itemize}
\end{proposition}

La seule propriété non triviale réside dans le fait que $\Lambda
\to\R\pi_\star\Lambda$ soit un monomorphisme scindé. En remarquant
que la projection $\Drap(\mathscr E)\to X$ peut s'écrire comme un composé
de $r$ projections de fibrés projectifs, ceci se déduit de la formule du
fibré projectif\;\footnote{Plus précisément, Grothendieck a montré
(cf.~\cite{PFTI@Grothendieck}, ou \sga{6}{vi}{4.6} pour le même argument
dans le cas de la $K$-théorie algébrique) que la théorie des classes de
Chern permettait de calculer l'algèbre de cohomologie des fibrés de
drapeaux, fussent-ils incomplets.}.

\begin{proposition}[Principe de scindage II]\label{16-proposition-scindage-2}
Soit $X$ un $\ZZn$-schéma. Soit $(E):0\to \mathscr E' \to \mathscr E \vers p 
\mathscr E''\to 0$
une suite exacte courte de fibrés vectoriels sur $X$. On note $\Sect(E)$
le $X$-schéma défini par le fait que pour tout $X$-schéma $f\colon Y\to
X$, l'ensemble des $X$-morphismes $Y\to \Sect(E)$ s'identifie
naturellement à l'ensemble des sections de la surjection de fibrés
vectoriels 
$f^\star(p)\colon f^\star \mathscr E\to f^\star \mathscr E''$ sur
$Y$\;\footnote{Je remercie Dennis Eriksson de m'avoir signalé cette
construction.}.
Le $Y$-schéma $\Sect(E)$ est naturellement muni d'une structure de 
torseur 
sous le $Y$-schéma en groupes vectoriel d'homomorphismes 
$\SheafHom(\mathscr E'',\mathscr E')$.
Notons $\pi\colon \Sect(E)\to X$ la projection. Les propriétés
suivantes sont satisfaites :
\begin{itemize}
\item l'image inverse par $\pi\colon \Sect(E)\to X$
de la suite exacte de fibrés vectoriels
$E$ est (canoniquement) scindée ;
\item le morphisme canonique $\Lambda\to \R\pi_\star\Lambda$
est un isomorphisme dans $\detaleplus X$.
\end{itemize}
\end{proposition}

L'existence de $\Sect(E)$ est évidente, la question étant de
nature locale sur $X$. Localement pour la topologie de Zariski sur $X$, la
projection $\pi$ est la projection depuis un espace affine, l'isomorphisme
$\Lambda\iso \R\pi_\star\Lambda$ résulte alors de l'invariance par
homotopie de la cohomologie étale pour les $\ZZn$-schémas (\sga{4}{xv}{2.2}).

\medskip

Démontrons le théorème~\ref{16-theoreme-classes-de-chern}. Grâce aux
propositions~\ref{16-proposition-scindage-1} et
\ref{16-proposition-scindage-2}, l'unicité est évidente. Il s'agit donc de
construire une théorie des classes de Chern satisfaisant les propriétés
demandées. Soit $\mathscr E$ un fibré vectoriel (que l'on peut
supposer de rang constant $r$) sur un $\ZZn$-schéma $X$. On considère le
fibré projectif $\PP(\mathscr E)$ sur $X$.
On note $\xi=c_1(\OO(1))$. D'après la
formule du fibré projectif
(théorème~\ref{16-theoreme-formule-du-fibre-projectif}), il existe d'uniques éléments, notés
$c_i(\mathscr E)\in \HH^{2i}(X,\Lambda(i))$ pour $1\leq i\leq r$ tels 
que l'on ait la
relation
\[\xi^r - c_1(\mathscr E)\xi^{r-1}+c_2(\mathscr E)\xi^{r-2}
+\dots+(-1)^r c_r(\mathscr E)=0\in
\HH^{2r}(\PP(\mathscr E),\Lambda(r))\;\text{.}\]
On pose $c_0(\mathscr E)=1$ et $c_i(\mathscr E)=0$ pour $i>r$.
Dans le cas où $\mathscr E$ est un
fibré en droites, $\PP(\mathscr E)\simeq X$
et $\OO(1)\simeq \mathscr E$, ce qui 
montre que cette définition étend la précédente pour les fibrés en droites.
La seule propriété non évidente est la formule de
Cartan-Whitney. Par principe de scindage
(propositions~\ref{16-proposition-scindage-1} et \ref{16-proposition-scindage-2}),
il suffit d'établir la formule
suivante :

\begin{lemme}
Soit $X$ un $\ZZn$-schéma. Soit $(\mathscr L_i)_{1\leq i\leq r}$ 
une famille finie
de fibrés en droites sur $X$, soit $\mathscr E=\bigoplus_{1\leq i\leq r}
\mathscr L_i$ leur
somme directe. Dans $\HH^{2r}(\PP(\mathscr E),\Lambda(r))$, on a la relation :
\[\prod_{i=1}^r(\xi-c_1(\mathscr L_i))=0\]
où $\xi=c_1(\OO(1))$. Autrement dit, 
\[c_t(\mathscr E)=\prod_{i=1}^r c_t(\mathscr L_i)\;\text{.}\]
\end{lemme}

L'argument qui suit est inspiré de \cite{PF@Panin-Smyrnov}. 
Pour $1\leq i\leq r$, on note $H_i$ l'hyperplan projectif de $\PP(\mathscr E)$
défini par l'inclusion $\mathscr L_i\to \mathscr E$. Notons $\pi\colon
\PP(\mathscr E)\to X$ la projection. Le morphisme canonique
$\pi^\star\mathscr L_i\to \OO(1)$ induit un isomorphisme sur l'ouvert
complémentaire de $H_i$ dans $\PP(\mathscr E)$.
On en déduit que l'élément $\xi-c_1(\mathscr L_i)$
de $\HH^2(X,\Lambda(1))$ peut être relevé en un élément $x_i$ du groupe de 
cohomologie à supports $\HH^2_{H_i}(X,\Lambda(1))$\;\footnote{Pour 
le moment,
peu importe de fixer un relèvement canonique.}.
Le produit des éléments $x_i$ vit naturellement dans le
groupe de cohomologie à support
$\HH^{2r}_{\cap_{1\leq i \leq r} H_i}(\PP(E),\Lambda(i))$ qui est nul puisque
l'intersection de ces $r$ hyperplans est vide ; on en
déduit la formule voulue par oubli du support.

\begin{proposition}
Soit $\mathscr E$ un fibré vectoriel sur un $\ZZn$-schéma $X$.
Pour tout entier
naturel $i$, on a l'égalité :
\[c_i(\mathscr E^\dual)=(-1)^i c_i(\mathscr E)\;\text{;}\]
autrement dit, on a une formule de changement de variables :
\[c_t(\mathscr E^\dual)=c_{-t}(\mathscr E)\;\text{.}\]
\end{proposition}

Grâce à la relation de Cartan-Whitney et au principe de scindage, on peut 
se ramener au cas où $\mathscr E$ est un fibré en droites. Cela résulte alors 
du fait que $c_1\colon \Pic(X)\to \HH^2(X,\Lambda(1))$ soit un homomorphisme
de groupes.

%>>>1
\section{Morphismes de Gysin}%<<<1
\label{16-section-gysin}


Étant donné un morphisme d'intersection complète $X\vers f S$ entre
$\ZZn$-schémas vérifiant certaines hypothèses techniques, on va construire
un morphisme de Gysin $\Cl f\colon \Lambda\to f^?\Lambda$ où
$f^?=f^!(-d)[-2d]$ ($d$ est la dimension relative virtuelle de $f$).
Ces morphismes de Gysin seront
compatibles à la composition des morphismes d'intersection complète.

L'essentiel de cette section est consacrée à la construction de ces
morphismes de Gysin dans le cas des immersions régulières. Le morphisme trace
permettra de faire la construction dans le cas des morphismes lisses.
Ces deux définitions se recolleront pour donner la
définition~\ref{16-definition-cl-f-lissifiable} dans le cas général et le
théorème~\ref{16-theoreme-cl-f-transitivite} établira la compatibilité à la
composition de ces morphismes de Gysin.

\subsection{Première classe de Chern d'un pseudo-diviseur}%<<<2

Soit $\mathscr L$ un fibré en droites sur un $\ZZn$-schéma $X$, $Z$ un
fermé de $X$ et $U$ l'ouvert complémentaire. On suppose donnée une section
inversible $s\colon \OO_U\iso {\mathscr L}_{|U}$. Au couple $(\mathscr
L,s)$ est canoniquement associée une classe $c_1(\mathscr L,s)\in
\HH^2_Z(X,\Lambda(1))$ induisant $c_1(\mathscr L)\in \HH^2(X,\Lambda(1))$ par
oubli du support (construire un élément de $\HH^1_Z(X,\Gm)$ et
utiliser la suite exacte de Kummer).

La classe $c_1(\mathscr L,s)$ définit un morphisme
$\Lambda_Z=\Lambda_X/\Lambda_U\to \Lambda_X(1)\cro 2$ dans $\detaleplus X$. En
\guil{composant} un tel morphisme avec une classe de cohomologie de $Z$
représentée par un morphisme $\Lambda_Z\to\Lambda_Z(q)\cro p$, il vient que
$c_1(\mathscr L,s)$ induit des morphismes \[\Gys_{(\mathscr L,s)}\colon
\HH^p(Z,\Lambda(q))\to \HH^{p+2}_Z(X,\Lambda(q+1))\;\text{.}\]

\begin{definition2}
Si $Z\to X$ est une immersion régulière de codimension $1$ définie par un
Idéal (inversible) $\II$, on pose $\Gys_{Z\subset X}=\Gys_{(\II
,1_{X-Z})}$.
\end{definition2}

On prendra garde au fait que \emph{via} les identifications usuelles, le
faisceau inversible $\II$ et le diviseur effectif $Z$ ont des
classes opposées dans le groupe de Picard de $X$.


%>>>2
\subsection{Classes fondamentales généralisées}%<<<2

Pour étudier la compatibilité à la composition des classes
fondamentales définies dans \cite[§1]{AZUMINO@Fujiwara} dans le cas des
immersions régulières (cf.~\sga{6}{vii}{1.4}),
Ofer Gabber définit une classe fondamentale
généralisée pour une immersion fermée $Y\to X$ définie par un Idéal de
type fini $\II$. Cette construction n'est plus limitée aux
immersions régulières et est compatible aux 
changements de bases arbitraires, mais elle 
dépend d'une donnée supplémentaire, à savoir celle d'un fibré vectoriel
sur $Y$ se surjectant sur le faisceau conormal $\Normal Y X=\II/\II^2$.

\subsubsection{Éclatement modifié}%<<<3

Soit $Y\to X$ une immersion fermée entre $\ZZn$-schémas définie 
par un Idéal de type fini $\II$. On note $U$ l'ouvert complémentaire.
Soit $\mathscr E
\to \Normal Y X$ un épimorphisme de Modules sur $Y$ où 
$\mathscr E$ est un Module localement libre de rang fini.
On définit une $\mathscr O_X$-Algèbre graduée quasi-cohérente $\mathscr
A_\star$ par produit fibré de façon à 
avoir un carré cartésien de $\mathscr O_X$-Modules,
pour tout entier naturel $n$ :
\[
\xymatrix{
\ar[d]{\mathscr A}_n\ar[r] & \II^n \ar[d]\\
\SS^n\mathscr E\ar[r] & \II^n/\II^{n+1}}
\]
où l'algèbre symétrique $\SS^\star\mathscr E$ est prise sur le faisceau
d'anneaux $\mathscr O_Y=\mathscr O_X/\II$.

\begin{definition3}
On pose $\Ecl_{Y,\mathscr E}(X)=\Proj({\mathscr A}_\star)$ et on note
$\pi\colon \Ecl_{Y,\mathscr E}(X)\to X$ la projection.
\end{definition3}

\begin{remarque3}
Si $Y\to X$ est une immersion fermée
régulière et que $\mathscr E\to \Normal Y X$
est un isomorphisme, $\Ecl_{Y,\mathscr E}(X)$ s'identifie à l'éclaté de $Y$
dans $X$. C'est ce cas particulier que l'on généralise ici en vue d'obtenir
une construction compatible aux changements de base.
\end{remarque3}

\begin{proposition3}\label{16-proposition-proprietes-eclatement-modifie}
L'Algèbre ${\mathscr A}_0$ est isomorphe à ${\mathscr O}_X$, les Modules
${\mathscr A}_n$ sont de type fini, l'Algèbre graduée
${\mathscr A}_\star$ est engendrée par ${\mathscr A}_1$ et on a un 
isomorphisme canonique de $\mathscr O_Y$-Algèbres graduées 
$\mathscr A_\star \otimes_{\mathscr O_X}{\mathscr O_X/\II}\iso
\SS^\star \mathscr E$.
\end{proposition3}


L'assertion concernant $\mathscr A_0$ est tautologique.
Soit $n$ un entier naturel. Comme $\II^n \to \II^n/\II^{n+1}$ est un
épimorphisme, la projection $\mathscr A_n\to \SS^n\mathscr E$ est aussi un
épimorphisme et si on note $\mathscr K_n$ son noyau, on a un isomorphisme
$\mathscr K_n\iso \II^{n+1}$. Par dévissage, il en résulte que $\mathscr
A_n$ est un Module de type fini.

Puisque $\SS^n\mathscr E$ est un $\mathscr O_X/\II$-Module, $\mathscr K_n$
contient $\II\cdot\mathscr A_n$. Comme $\mathscr E\to \II/\II^2$ est un
épimorphisme, le morphisme $\SS^n\mathscr E\to \II^n/\II^{n+1}$ est aussi
un épimorphisme, ce qui implique que la projection $\mathscr A_n\to \II^n$
est un épimorphisme. Par conséquent, l'inclusion
$\II\cdot\mathscr A_n\to\II\cdot
\II^n=\II^{n+1}\simeq \mathscr K_n$ est un isomorphisme : $\mathscr
K_n=\II\cdot \mathscr A_n$. Ceci permet d'obtenir l'isomorphisme
$\mathscr A_\star\otimes_{\mathscr O_X} {\mathscr O_X/\II}\iso
\SS^\star \mathscr E$.

Pour montrer que le morphisme évident $\mathscr A_1^{\otimes n}\to
\mathscr A_n$ de Modules
est un épimorphisme, il suffit, d'après le lemme de
Nakayama, de le tester après passage aux corps résiduels de $X$.
Au-dessus de l'ouvert $U$, c'est évident ; au-dessus de $Y$, cela résulte
de l'isomorphisme 
$\mathscr A_\star\otimes_{\mathscr O_X} {\mathscr O_X/\II}\iso 
\SS^\star \mathscr E$.

\begin{corollaire3}\label{16-corollaire-structure-eclatement-modifie}
Le morphisme $\pi\colon \Ecl_{Y,\mathscr E}(X)\to X$ est projectif et on
dispose d'isomorphismes canoniques $\pi^{-1}(U)\iso U$ et
$\pi^{-1}(Y)\simeq \PP(\mathscr E)$.
\end{corollaire3}

L'isomorphisme au-dessus de $U$ est évident. Compte tenu de \ega{ii}{3.5.3},
celui décrivant $\pi^{-1}(Y)$
se déduit de l'isomorphisme de $\mathscr O_Y$-Algèbres graduées $\mathscr
A_\star \otimes_{\mathscr O_X}\mathscr O_Y\iso \SS^\star\mathscr E$.

\begin{proposition3}\label{16-proposition-cb-eclatements-modifies}
Soit $p\colon X'\to X$ un morphisme.
On pose $Y'=Y\times_X X'$ et $\mathscr E'=p^\star \mathscr E$. On dispose
d'un épimorphisme évident $\mathscr E'\to\Normal {Y'} {X'}$. Le
morphisme canonique
\[\Ecl_{Y',\mathscr E'}(X')\to \Ecl_{Y,\mathscr E}(X)\times_X X'\]
est une nil-immersion.
\end{proposition3}

Notons $\mathscr A'_\star$ la $\mathscr O_{X'}$-Algèbre 
graduée quasi-cohérente donnant
naissance à $\Ecl_{Y',\mathscr E'}(X')$. On dispose d'un morphisme évident 
$p^\star \mathscr A_\star\to\mathscr A'_\star$ de $\mathscr
O_{X'}$-Algèbres graduées quasi-cohérentes. Pour tout entier, le morphisme
$p^\star \mathscr A_n\to \mathscr A'_n$ est un morphisme entre $\mathscr
O_{X'}$-Modules de type fini ; pour montrer qu'il s'agit d'un épimorphisme,
d'après le lemme de Nakayama, il suffit de vérifier que ce morphisme
induit un isomorphisme d'une part au-dessus de $U'=X'-Y'$ (c'est
évident) et d'autre part modulo l'idéal $\II'$ définissant $Y'$ dans $X'$
(cela résulte de la description donnée dans la 
proposition~\ref{16-proposition-proprietes-eclatement-modifie}). Le morphisme 
\[\Ecl_{Y',\mathscr E'}(X')\to \Ecl_{Y,\mathscr E}(X)\times_X X'\]
s'identifie au $X'$-morphisme évident
$\Proj(\mathscr A'_\star)\to \Proj (p^\star \mathscr A_\star)$
(\ega{ii}{3.5.3}) ; d'après ce
qui précède, il s'agit d'une immersion fermée. Le fait que ce morphisme
induise un isomorphisme au-dessus de $p^{-1}(U)$ et de $p^{-1}(Y)$
permet d'en déduire aussitôt que le morphisme induit au niveau des schémas
réduits associés
\[\Ecl_{Y',\mathscr E'}(X')_{\red}\to (\Ecl_{Y,\mathscr E}(X)\times_X
X')_\red\]
est un isomorphisme.

%>>>3
\subsubsection{Définition des classes}%<<<3

On se donne toujours une immersion fermée $i\colon Y\to X$ définie par un
Idéal $\II$ de type fini. On note $j\colon U\to X$ l'inclusion de l'ouvert
complémentaire ($\II$ étant de type fini, $j$ est un morphisme de 
type fini). On suppose donné un épimorphisme de $\mathscr O_Y$-Modules
$\mathscr E \to \Normal Y X$
avec $\mathscr E$ est localement libre de rang
fini. On note $\pi\colon \Ecl_{Y,\mathscr E}(X)\to X$ la projection de
l'éclatement modifié, $j'\colon U\to \Ecl_{Y,\mathscr E}(X)$ l'immersion
ouverte évidente et $i'\colon \PP(\mathscr E)\to \Ecl_{Y,\mathscr E}(X)$
l'immersion fermée donnée par le 
corollaire~\ref{16-corollaire-structure-eclatement-modifie}. On note $r$ le
rang du fibré vectoriel $\mathscr E$ que l'on suppose de rang constant pour
simplifier et on suppose $r>0$.

On a ainsi le diagramme suivant de schémas, où les carrés sont cartésiens :

\[
\xymatrix{
\ar[d]^{\pi'}\PP(\mathscr E)\ar[r]^{i'} & 
\Ecl_{Y,\mathscr E}(X)\ar[d]^\pi &\ar[l]_{j'} U\ar@{=}[d]\\
Y \ar[r]^i & X & \ar[l]_{i'} U}
\]

\begin{proposition3}
\label{16-proposition-cohomologie-eclatement}
Le morphisme évident $\Lambda\to \R\pi'_\star\Lambda$ dans $\detaleplus Y$
est un monomorphisme scindé : la formule du fibré projectif identifie son
conoyau à 
\[\bigoplus_{k=1}^{r-1} \Lambda(-k)\cro{-2k}\;\text{.}\]
Les morphismes évidents définissent un triangle distingué :
\[\Lambda\to \R\pi_\star\Lambda \to
i_\star\coker(\Lambda\to \R\pi'_\star \Lambda)\vers 0 \Lambda\cro 1\]
dans $\detaleplus X$. On peut le récrire sous la forme
\[\Lambda\to \R\pi_\star\Lambda \vers \rho
\bigoplus_{k=1}^{r-1}i_\star\Lambda(-k)\cro{-2k}\vers 0 \Lambda\cro
1\;\text{,}\]
le morphisme $\rho$ admettant une section canonique
donnée par les éléments
$c_1(\OO(1),1_U)^k$ de $\HH^{2k}_{\PP(\mathscr E)}(\Ecl_{Y,\mathscr
E}(X),\Lambda(k))$, identifiés à des morphismes 
$i_\star \Lambda(-k)\cro {-2k}\to \R\pi_\star \Lambda$ dans $\detaleplus X$.
\end{proposition3}

On note $L$ une résolution injective du faisceau constant $\Lambda$ vu
comme faisceau de $\Lambda$-modules sur le grand site étale
des schémas de type fini sur $X$. Pour tout morphisme de type fini $W\vers
p X$, on note $L_{|W}$ le complexe de faisceau de $\Lambda$-modules sur
$W_\et$ induit par $L$ ; on peut le voir comme un objet de $\detaleplus W$
isomorphe à $\Lambda$.

\begin{lemme3}
Le carré commutatif évident de complexes de faisceaux sur $X_{\et}$
est homotopiquement bicartésien :
\[
\xymatrix{
\ar[d]L_{|X} \ar[r] & i_\star L_{|Y} \ar[d]\\
\pi_\star L_{|\Ecl_{Y,\mathscr E}(X)}\ar[r] & \pi_\star i'_\star
L_{|\PP(\mathscr E)}}
\]
(ceci signifie par exemple
que le complexe simple associé à ce diagramme, identifié à
un complexe $3$-uple, est acyclique).
\end{lemme3}

Les complexes simples associés aux complexes doubles 
\[j_!L_{|U}\to L_{|X} \to i_\star L_{|Y}\]
et 
\[j'_!L_{|U}\to L_{|\Ecl_{Y,\mathscr E}(X)} \to i'_\star L_{|\PP(\mathscr E)}\]
de faisceaux sur $X$ et $\Ecl_{Y,\mathscr E}(X)$ respectivement sont
acycliques.
Choisissons un foncteur de résolution \guil{flasque} additif
$r$ sur la catégorie des faisceaux de $\Lambda$-modules sur 
${\Ecl_{Y,\mathscr E}(X)}_\et$ et notons abusivement
$\R\pi_\star$ le foncteur (additif)
de la catégorie des
complexes (bornés inférieurement) de faisceaux de $\Lambda$-modules
sur $\Ecl_{Y,\mathscr E}(X)$ vers la catégorie des
complexes de faisceaux de $\Lambda$-modules
sur $X$ défini par la formule $\R\pi_\star K = \Tot(\pi_\star r K)$,
ce foncteur préserve les
quasi-isomorphismes et induit
le foncteur $\R\pi_\star\colon \detaleplus {\Ecl_{Y,\mathscr E}(X)}\to
\detaleplus X$ usuel.

On obtient ainsi un diagramme commutatif de complexes de faisceaux de
$\Lambda$-modules sur $X$ :
\[
\xymatrix{
j_!L_{|U} \ar[r]\ar[d] & L_{|X} \ar[r] \ar[d]& i_\star L_{|Y} \ar[d]\\
\R\pi_\star j'_! L_{|U} \ar[r] & \R\pi_\star
L_{|\Ecl_{Y,\mathscr E}(X)}\ar[r] & \R\pi_\star i'_\star
L_{|\PP(\mathscr E)}}
\]
Les lignes de ce diagramme constituent des complexes doubles dont les
complexes simples associés sont acycliques. D'après le théorème de
changement de base pour un morphisme propre, le morphisme 
$j_!L_{|U}\to \R\pi_\star j'_! L_{|U}$ est un quasi-isomorphisme.
On en déduit que le carré de droite est homotopiquement bicartésien, ce qui
permet de conclure.

\medskip

Revenons à la démonstration de la
proposition~\ref{16-proposition-cohomologie-eclatement},
la formule du fibré projectif pour $\PP(\mathscr E)$ implique que l'on a un
triangle distingué dans $\detaleplus X$ :
\[i_\star \Lambda\to \R\pi_\star i'_\star \Lambda
\to \bigoplus_{i=1}^{r-1}i_\star \Lambda(-i)\cro{-2i} \vers 0 i_\star
\Lambda\cro 1\;\text{.}\]
En considérant les colonnes du carré homotopiquement bicartésien donné par
le lemme, on peut conclure à l'existence d'un triangle distingué
\[
\xymatrix{
\Lambda \to \R\pi_\star \Lambda \to 
\displaystyle\bigoplus_{i=1}^{r-1} i_\star
\Lambda(-i)\cro{-2i}\to\Lambda\cro 1\;\text{.}}
\]
Ce triangle est scindé par les puissances de l'élément $c_1(\OO(1),1_U)$ ; le
morphisme de droite est donc nul, ce qui achève la démonstration de la
proposition.

\begin{corollaire3}\label{16-corollaire-caracterisation-image-restriction}
La suite suivante, dont les morphismes sont évidents, est exacte :
\[0\to \HH^{2r}_Y(X,\Lambda(r)) \to 
\HH^{2r}_{\PP(\mathscr E)}(\Ecl_{Y,\mathscr E}(X),\Lambda(r))\to
\coker(\HH^{2r}(Y,\Lambda(r))\to \HH^{2r}(\PP(\mathscr E),\Lambda(r)))\to 0\]
\end{corollaire3}

L'énoncé de ce corollaire serait bien évidemment juste en tout bidegré
$(p,q)$ au lieu de $(2r,r)$, mais nous n'utiliserons que ce cas
particulier.

\bigskip

On note $\Gys\colon \HH^p(\PP(\mathscr E),\Lambda(q))\to
\HH^{p+2}_{\PP(\mathscr E)}(\Ecl_{Y,\mathscr E}(X),\Lambda(q+1))$ le
morphisme de Gysin associé au pseudo-diviseur $(\OO(1),1_U)$ sur 
$\Ecl_{Y,\mathscr E}(X)$ et $\xi=c_1(\OO(1))\in \HH^2(\PP(\mathscr
E),\Lambda(1))$. Le lemme suivant est évident :

\begin{lemme3}\label{16-lemme-res-gys-egale-xi}
Le morphisme composé
\[\HH^p(\PP(\mathscr E),\Lambda(q))\vers {\Gys} \HH^{p+2}_{\PP(\mathscr
E)}(\Ecl_{Y,\mathscr E}(X),\Lambda(q+1))\to  \HH^{p+2}(\PP(\mathscr
E),\Lambda(q+1))\;\text{,}\]
où la flèche de droite est le morphisme de restriction, est la
multiplication par $\xi$.
\end{lemme3}

\begin{definition3}
On définit un élément $\Clf {i,\mathscr E}$ de $\HH^{2r-2}(\PP(\mathscr
E),\Lambda(r-1))$ par la formule :
\[\Clf {i,\mathscr E}=-\cro{\xi^{r-1}-c_1(\mathscr
E)\xi^{r-2}+\dots+ (-1)^{r-1}c_{r-1}(\mathscr E)}\;\text{.}\]
\end{definition3}

Le lemme suivant résulte aussitôt de la construction des classes de Chern :

\begin{lemme3}\label{16-lemme-xi-clf}
Dans $\HH^{2r}(\PP(\mathscr E),\Lambda(r))$, on a l'égalité
\[\xi\Clf{i,\mathscr E}=(-1)^rc_r({\mathscr E})\;\text{.}\]
\end{lemme3}

\begin{definition3}
Compte tenu du corollaire~\ref{16-corollaire-caracterisation-image-restriction}, 
les lemmes~\ref{16-lemme-res-gys-egale-xi}
et \ref{16-lemme-xi-clf} 
montrent que l'élément
$\Gys(\Clf{i,\mathscr E})\in \HH^{2r}_{\PP(\mathscr E)}(\Ecl_{Y,\mathscr
E}(X),\Lambda(r))$ provient par restriction d'un unique élément de
$\HH^{2r}_Y(X,\Lambda(r))$, noté $\Cl{i,\mathscr E}$.
\end{definition3}

%>>>3
\subsubsection{Propriétés des classes généralisées}%<<<3

\begin{proposition3}
\label{16-proposition-classes-generalisees-cb}
La formation des classes généralisées $\Cl{i,\mathscr E}$ 
et $\Clf {i,\mathscr E}$ est compatible à tout changement de base $X'\to X$.
\end{proposition3}

Compte tenu de la 
proposition~\ref{16-proposition-cb-eclatements-modifies}, ceci résulte
aussitôt des définitions.

\begin{proposition3}
\label{16-proposition-classes-generalisees-changement-de-fibre}
Soit ${\mathscr E}'\to {\mathscr E}$ un épimorphisme
de Modules localement libres sur $Y$. Soit $\mathscr K$ le noyau de cet
épimorphisme. On suppose que ${\mathscr E}'$ est de rang constant $r'$.
On a alors la relation
\[\Cl {i,\mathscr E'}= (-1)^{r'-r}c_{r'-r}({\mathscr K})\cdot\Cl
{i,\mathscr E}\]
dans $\HH^{2r'}_Y(X,\Lambda(r'))$ où on a utilisé les accouplements
canoniques \[\HH^a(Y,\Lambda(b))\otimes 
\HH^{a'}_Y(X,\Lambda(b'))\to  \HH^{a+a'}_Y(X,\Lambda(b+b'))\;\text{.}\]
\end{proposition3}

On dispose d'une immersion fermée de 
$\Ecl_{Y,\mathscr E}(X)$ dans 
$\Ecl_{Y,{\mathscr E}'}(X)$, ce qui permet de considérer la composition
suivante de flèches de restriction :
\[\HH^{2r'}_Y(X,\Lambda(r'))\to \HH^{2r'}_{\PP(\mathscr E')}(\Ecl_{Y,\mathscr
E'}(X),\Lambda(r'))\to \HH^{2r'}_{\PP(\mathscr E)}(\Ecl_{Y,\mathscr
E}(X),\Lambda(r'))\;\text{.}\]
Cette composée étant injective, il s'agit de montrer que les images des
deux éléments considérés dans $\HH^{2r'}_{\PP(\mathscr E)}(\Ecl_{Y,\mathscr
E}(X),\Lambda(r'))$
sont égales, mais comme ces deux éléments sont naturellement définis
comme étant des images d'éléments de $\HH^{2r'-2}(\PP(\mathscr
E),\Lambda(r'-1))$ par le morphisme $\Gys$ associé au fibré en droites
$\OO(1)$ sur $\Ecl_{Y,\mathscr E}(X)$ canoniquement trivialisé
sur $X-Y$, on se ramène à montrer l'égalité
\[{\Clf{i,{\mathscr E}'}}_{|\PP(\mathscr E)}=
(-1)^{r'-r}c_{r'-r}({\mathscr K})\cdot\Clf{i,\mathscr
E}\]
dans $\HH^{2r'}(\PP(\mathscr E),\Lambda(r'))$. Ceci ne fait plus intervenir
que le schéma $Y$ et l'épimorphisme ${\mathscr E}'\to \mathscr E$ de fibrés
vectoriels sur $Y$. Il s'agit d'une identité \guil{universelle} dans cette
situation ; on peut appliquer le principe de scindage et faire une
récurrence sur la différence $r'-r$ pour se ramener au cas où $r'=r+1$,
c'est-à-dire que $\mathscr K$ est un fibré en droites.
Les relations de Cartan-Whitney entre les classes de Chern de
$\mathscr E$, ${\mathscr E}'$ et $\mathscr K$ donnent aussitôt l'égalité
\[{\Clf{i,\mathscr E'}}_{|\PP(\mathscr E)} = \cro{\xi \Clf{i,\mathscr
E}-(-1)^rc_r({\mathscr E})}-c_1(\mathscr K)\cdot\Clf{i,\mathscr E}\;\text{,}\]
ce qui donne bien la relation voulue puisque $\xi\Clf{i,\mathscr
E}=(-1)^rc_r({\mathscr E})$ (cf.~lemme~\ref{16-lemme-xi-clf}).

%>>>3
%>>>2
\subsection{Immersions régulières}%<<<2

On rappelle que la notion d'immersion régulière est définie dans
\sga{6}{vii}{1.4}.

\begin{definition2}\label{16-definition-cl-i}
Soit $i\colon Y\to X$ une immersion régulière entre $\ZZn$-schémas.
On pose $i^?=i^\star(c)\cro{2c}\colon \detaleplus X\to \detaleplus Y$ où $c$ est la
codimension de $i$.
On définit un morphisme 
$\Cl i \colon \Lambda\to i^?\Lambda$ dans $\detaleplus Y$ de la façon suivante.
Quitte à décomposer $Y$ en réunion disjointe d'ouverts-fermés, on peut
supposer que la codimension $c$ de $i$ est constante.
Si $c=0$, $i$ est l'inclusion d'un ouvert, $\Cl i$ est
l'isomorphisme évident. Dans le cas où $c>0$,
choisissons un ouvert $U$ de $X$ dans lequel $Y$ est un sous-schéma fermé,
notons $i'\colon Y\to U$ cette immersion fermée. Le faisceau conormal
$\Normal Y X$ de $Y$ dans $X$ est un fibré vectoriel de rang $c$ sur $Y$
muni de l'épimorphisme tautologique $\Normal Y X\to\Normal Y X$ ; on peut
donc considérer la classe $\Cl {i'}=\Cl {i',\Normal Y X}\in
\HH^{2c}_Y(U,\Lambda(c))$, que l'on identifie à un morphisme $\Cl i\colon
\Lambda\to
i'^?\Lambda\simeq i^?\Lambda$ dans $\detaleplus Y$ ; il est évident que la
construction ne dépend pas de l'ouvert intermédiaire $U$.
\end{definition2}

Le théorème suivant généralise l'énoncé établi dans
\cite[proposition~1.2.1]{AZUMINO@Fujiwara} :

\begin{theoreme2}\label{16-theoreme-composition-cl-immersions-regulieres}
Si $Z\vers i Y$ et $Y \vers j X$ sont deux 
immersions régulières composables, le
diagramme suivant est commutatif dans $\detaleplus Z$ :
\[
\xymatrix{
\ar[rd]_{\Cl {j\circ i}}\Lambda\ar[r]^{\Cl i}  &
   i^?\Lambda \ar[d]^{i^?(\Cl j)} \\
&   i^?j^?\Lambda}
\]
\end{theoreme2}

On peut évidemment supposer que les immersions $i$ et $j$ sont des
immersions fermées et que les codimensions de $i$ et de $j$ sont constantes,
de valeurs respectives $m$ et $n$. Si $m=0$ ou $n=0$, c'est trivial ; on
suppose donc que $m>0$ et $n>0$.

\begin{lemme2}
On peut supposer que $n=1$ (\emph{i.e.} $j$ est de codimension $1$).
\end{lemme2}

On éclate $Y$ dans $X$ pour obtenir le diagramme suivant où les carrés sont
cartésiens :
\[
\xymatrix{
P'\ar[r]^{i'} \ar[d]^{p'}& P\ar[r]^{j'}\ar[d]^p & \Ecl_Y(X)\ar[d]^\pi \\
Z\ar[r]^i& Y\ar[r]^j & X}
\]

L'idée de la démonstration est d'utiliser une sorte de formule d'excès
d'intersection (cf.~\cite[theorem~6.3]{IT@Fulton} pour une formulation dans la
théorie de Chow) pour les immersions $j\circ i\colon Z\to X$ et $j\colon Y
\to X$ relativement au changement de base $\pi\colon \Ecl_Y(X)
\to X$ qui va faire
chuter la codimension de ces immersions fermées régulières.

On a des isomorphismes canoniques $P=\PP(\Normal Y X)$ et $P'=\PP({\Normal
Y X}_{|Z})$.
On vérifie facilement que $P\to \Ecl_Y(X)$ est une 
immersion fermée régulière de
codimension $1$. Par changement de base lisse, $P'\to P$ est une immersion
fermée régulière de codimension $m$. On suppose que $i'^?(\Cl {j'})\circ
\Cl {i'}=\Cl{j'\circ i'}$ et on veut montrer que $i^?(\Cl {j})\circ
\Cl {i}=\Cl{j\circ i}$. Les morphismes à comparer s'identifient à des
éléments de $\HH^{2(m+n)}_Z(X,\Lambda(m+n))$ (on fera ce type
d'identifications jusqu'à la fin de la démonstration). La
proposition~\ref{16-proposition-cohomologie-eclatement} implique que
l'application
\[\pi^\star\colon \HH^{2(m+n)}_Z(X,\Lambda(m+n))\to
\HH^{2(m+n)}_{P'}(\Ecl_Y(X),\Lambda(m+n))\]
est injective, il suffit donc de comparer les classes après application de
$\pi^\star$.

Considérons $\pi^\star(\Cl {j\circ i})\in
\HH^{2(m+n)}_{P'}(\Ecl_Y(X),\Lambda(m+n))$. La classe $\Cl {j\circ i}$ est la
classe généralisée $\Cl {j\circ i,\Normal Z X}$, la 
proposition~\ref{16-proposition-classes-generalisees-cb} implique l'égalité
\[\pi^\star \Cl {j\circ i}=\Cl {j'\circ i',\pi^\star\Normal Z X}\]
où est sous-entendu l'épimorphisme de fibrés vectoriels $p'^\star \Normal
Z X\to \Normal {P'} {\Ecl_Y(X)}$ dont on note $\mathscr E'$ le noyau, qui est un
$\mathscr O_{P'}$-Module localement libre de rang $n-1$. La
proposition~\ref{16-proposition-classes-generalisees-changement-de-fibre}
donne alors l'égalité
\[\pi^\star\Cl{j\circ i}=c_{n-1}({\mathscr E'}^\dual)\cdot \Cl {j'\circ i'}\]
où l'on a utilisé l'accouplement
\[\HH^{2(n-1)}(P',\Lambda(n-1))\times \HH^{2(m+1)}_{P'}(\Ecl_Y(X),\Lambda(m+1))
\to \HH^{2(m+n)}_{P'}(\Ecl_Y(X),\Lambda(m+n))\;\text{.}\]
La composition des classes admise provisoirement pour les immersions $j'$
et $i'$ donne l'égalité
\[\Cl {j'\circ i'}=\Cl {i'}\cdot \Cl{j'}\]
\emph{via} l'accouplement
\[\HH^{2m}_{P'}(P,\Lambda(m))\times \HH^2_P(\Ecl_Y(X),\Lambda(1))\to
\HH^{2(m+1)}_{P'}(\Ecl_Y(X),\Lambda(n+1))\;\text{.}\]
On a ainsi obtenu :
\[\pi^\star \Cl{j\circ i}=c_{n-1}({\mathscr E'}^\dual)\cdot \Cl {i'}\cdot \Cl
{j'}\;\text{.}\]
Notons $\mathscr E$ le noyau de l'épimorphisme $p^\star \Normal Y
X\to\Normal {P'} {\Ecl_Y(X)}$. Il vient aussitôt que dans le diagramme évident
de Modules sur $P'$ qui suit, les lignes et les colonnes sont exactes :
\[
\xymatrix{
& 0 \ar[d] & 0 \ar[d] &   & \\
0 \ar[r] & \ar[d]i'^\star\mathscr E\ar[r] & \ar[d]\mathscr E'\ar[r] & 0
\ar[d] & \\
0 \ar[r] & \ar[d]i'^\star p^\star \Normal Y X \ar[r] & \ar[d]
p'^\star \Normal Z X\ar[r]
&\ar[d] p'^\star \Normal Z Y \ar[r] & 0 \\
0 \ar[r] & \ar[d]i'^\star \Normal {P} {\Ecl_Y(X)} \ar[r] & \ar[d]
\Normal {P'} {\Ecl_Y(X)}\ar[r] & \ar[d] \Normal {P'} P \ar[r] & 0 \\
&0 &0 &0 &}
\]
En particulier, on obtient un isomorphisme canonique $i'^\star \mathscr
E\iso \mathscr E'$, d'où $i'^\star c_{n-1}({\mathscr E}^\dual)
=c_{n-1}({\mathscr E'}^\dual)\in \HH^{2(n-1)}(P',\Lambda(n-1))$. On en déduit :
\[\pi^\star \Cl{j\circ i}=c_{n-1}({\mathscr E'}^\dual)
\cdot \Cl {i'}\cdot \Cl{j'}=
\Cl {i'}\cdot c_{n-1}({\mathscr E}^\dual)\cdot \Cl {j'}\;\text{.}\]
On utilise implicitement dans ces notations l'associativité des structures
multiplicatives permettant par exemple de définir une application
\[\HH^{2m}_{P'}(P,\Lambda(m))\times \HH^{2(n-1)}(P,\Lambda(n-1))\times
\HH^2_P(\Ecl_Y(X),\Lambda(1))\to \HH^{2(m+n)}(\Ecl_Y(X),\Lambda(m+n))\]
sans qu'il y ait à s'inquiéter de l'ordre dans lequel les multiplications
sont faites.
Les propositions~\ref{16-proposition-classes-generalisees-changement-de-fibre}
et \ref{16-proposition-classes-generalisees-cb} impliquent les égalités
suivantes :
\[c_{n-1}({\mathscr E}^\dual)\cdot \Cl{j'}=\Cl{j',p^\star \Normal Y X}
=\pi^\star \Cl{j}\in \HH^{2n}_P(\Ecl_Y(X),\Lambda(n))\;\text{.}\]
Le morphisme $p$ étant lisse, on a aussitôt $\Cl {i'}=\pi^\star \Cl i$. On
a ainsi obtenu l'égalité voulue :
\[\pi^\star \Cl{j\circ i}=\pi^\star \Cl i\cdot \pi^\star \Cl j\;\text{,}\]
ce qui achève la démonstration du lemme.

\bigskip

On est ramené à établir le
théorème~\ref{16-theoreme-composition-cl-immersions-regulieres} dans le cas où
$j$ est de codimension $1$. On pose maintenant $P=\PP(\Normal Z X)$ et
$P'=\PP(\Normal Z Y)$. Le diagramme suivant récapitule la situation :

\[
\xymatrix{
P\ar[r] & \pi^{-1}(Y)\ar[r] & \Ecl_Z(X) \ar[dd]^\pi\\
P'\ar[u]\ar[r]\ar[d] & \ar[d]\Ecl_Z(Y)\ar[u] & \\
Z \ar[r]^i & Y \ar[r]^j & X}
\]

On veut établir l'égalité suivante dans $\HH^{2m+2}_Z(X,\Lambda(m+1))$ :
\[\Cl {j\circ i}=\Cl i\cdot \Cl j\;\text{.}\]
D'après la proposition~\ref{16-proposition-cohomologie-eclatement}, il suffit
de vérifier cette égalité dans $\HH^{2m+2}_{P}(\Ecl_Z(X),\Lambda(m+1))$ après
application de $\pi^\star$.

Par définition, la classe $\Cl{j\circ i}\in \HH^{2m+2}_Z(X,\Lambda(m+1))$ se
\guil{restreint} en un élément
\[\gamma=\Gys_{P\subset \Ecl_Z(X)}(\Clf{j\circ i})\]
dans $\HH^{2m+2}_P(\Ecl_Z(X),\Lambda(m+1))$ où $\Clf{j\circ i}\in
\HH^{2m}(P,\Lambda(m))$.

Notons $\II$ l'Idéal de $Y$ dans $X$, ${\II}_P$ celui de $P$
dans $\Ecl_Z(X)$ et $\tilde{\II}$ celui de $\Ecl_Z(Y)$ dans $\Ecl_Z(X)$.
On a un isomorphisme canonique de faisceaux inversibles sur $\Ecl_Z(X)$ :
\[\pi^\star\II\simeq {\II}_P\otimes \tilde{\mathscr
I}\;\text{.}\]
Cet isomorphisme est compatible aux trivialisations données sur
$\pi^{-1}(V)$ où $V=X-Y$. On obtient ainsi une égalité dans le 
groupe des classes d'équivalences de tels pseudo-diviseurs, ce qui 
permet de décomposer $c_1(\pi^\star\II, 1_{\pi^{-1}(V)})\in
\HH^2_{\pi^{-1}(V)}(\Ecl_Z(X),\Lambda(1))$
en une somme de deux composantes :
\[\pi^\star(c_1(\II,1_V))=c_1(\pi^\star\II, 1_{\pi^{-1}(V)})
= c_1({\II}_P,1_{\Ecl_Z(X)-P})+c_1(\tilde{\mathscr
I},1_{\Ecl_Z(X)-\Ecl_Z(Y)})\;\text{.}\]

On en déduit une décomposition
\[\pi^\star\Cl i\cdot\pi^\star \Cl j=\alpha+\beta\] dans
$\HH^{2m+2}_P(\Ecl_Z(X),\Lambda(m+1))$ où
\begin{eqnarray*}
\alpha &=&-\Gys_{P\subset \Ecl_Z(X)}({\Cl i}_{|P})\;\text{,} \\
\beta &=& -\Gys_{\Ecl_Z(Y)\subset \Ecl_Z(X)}(\Gys_{P'\subset \Ecl_Z(Y)}(\Clf i))\;\text{.}
\end{eqnarray*}
Le calcul de $\Cl k$ où $k$ est l'inclusion de l'intersection de diviseurs
de Cartier s'intersectant transversalement dans le schéma ambiant réalisé dans 
\cite[proposition~1.1.4]{AZUMINO@Fujiwara}
permet d'obtenir l'égalité d'opérateurs suivante :
\[\Gys_{\Ecl_Z(Y)\subset \Ecl_Z(X)}\circ \Gys_{P'\subset \Ecl_Z(Y)} = 
\Gys_{P\subset \Ecl_Z(X)}\circ \Gys_{P'\subset P}\;\text{.}\]
Notre but est d'établir l'égalité $\gamma=\alpha+\beta$. Les calculs
précédents permettent d'écrire chacun des éléments $\alpha$, $\beta$ et
$\gamma$ comme des images par le morphisme $\Gys_{P\subset \Ecl_Z(X)}$
de classes
$\tilde{\alpha}$, $\tilde{\beta}$ et $\tilde{\gamma}$ dans
$\HH^{2m}(P,\Lambda(m))$ :
\begin{eqnarray*}
\tilde{\alpha} &=& -{\Cl i}_{|P}\;\text{,}\\
\tilde{\beta} &=& -\Gys_{P'\subset P}(\Clf i)\;\text{,}\\
\tilde{\gamma} &=& \Clf {j\circ i}\;\text{.}
\end{eqnarray*}
On est ainsi ramené à établir l'égalité
$\tilde{\gamma}=\tilde{\alpha}+\tilde{\beta}$
dans $\HH^{2m}(P,\Lambda(m))$.

D'après les lemmes~\ref{16-lemme-res-gys-egale-xi} et \ref{16-lemme-xi-clf}, on a 
${\Cl i}_{|Z}=(-1)^{m}c_m(\Normal Z Y)$. On en déduit l'égalité
\[\tilde{\alpha}=(-1)^{m+1}c_m(\Normal Z Y)\;\text{.}\]
Pour calculer $\tilde{\beta}$, on observe que l'Idéal de $P'$ dans $P$
s'identifie au faisceau inversible $\mathscr K
\otimes_{{\mathscr O}_Z} \OO(-1)$
où $\mathscr K={\Normal Y X}_{|Z}$ est le noyau de l'épimorphisme
$\Normal Z X\to \Normal Z Y$. On en déduit 
\[\tilde{\beta}= (c_1(\mathscr K)-\xi)\cdot\cro{\xi^{m-1}-c_1(\Normal Z
Y)\xi^{m-2}+\dots + (-1)^{m-1}c_{m-1}(\Normal Z Y)}\;\text{.}\]
Par ailleurs, la définition de $\tilde{\gamma}$ donne l'égalité :
\[\tilde{\gamma} = -\cro{\xi^m-c_1(\Normal Z X)\xi^{m-1}+(-1)^mc_m(\Normal
Z X)}\;\text{.}\]
La formule de Cartan-Whitney appliquée à la suite exacte courte
\[0\to \mathscr K\to \Normal Z X\to\Normal Z Y\to 0\]
de fibrés vectoriels sur $Z$ permet d'obtenir aussitôt la relation voulue
$\tilde{\gamma}=\tilde{\alpha}+\tilde{\beta}$, ce qui achève la
démonstration du théorème.

\bigskip

La classe que l'on a définie est évidemment compatible avec celle de
\sgaquatreetdemi{Cycle}{2.2} :

\begin{proposition2}\label{16-proposition-compatibilite-sga-quatre-et-demi}
Soit $i\colon Y\to X$ une immersion régulière de codimension $c$
entre $\ZZn$-schémas.
Le morphisme de faisceaux $\Lambda\to {\mathscr H}^{2c}(i^!\Lambda(c))$
induit par le morphisme $\Cl i\colon \Lambda\to i^?\Lambda$ 
est donné par la classe $\clsga Y$ de \sgaquatreetdemi{Cycle}{2.2}.
\end{proposition2}

%>>>2
\subsection{Morphismes lisses}%<<<2

Soit $p\colon X\to S$ un morphisme lisse compactifiable
de $\ZZn$-schémas de dimension
relative $d$. D'après
\sga{4}{xviii}{2.9}, on dispose d'un morphisme trace
\[\Tr_p\colon \R^{2d}p_!\Lambda(d)\to \Lambda\;\text{,}\]
que l'on peut réinterpréter sous la forme d'un morphisme
\[\R p_!\Lambda(d)\cro{2d}\to \Lambda\]
dans $\detaleplus S$ (en effet, d'après le théorème de changement de base pour un
morphisme propre et \sga{4}{x}{4.3}, les faisceaux $\R^ip_!\Lambda$
sont nuls pour $i>2d$). 

\begin{definition2}\label{16-definition-cl-p}
Soit $p\colon X\to S$ un morphisme lisse compactifiable
de $\ZZn$-schémas. Le morphisme 
$\Cl p\colon \Lambda\to p^?\Lambda$\;\footnote{On rappelle que l'on a posé
$p^?=p^!(-d)\cro{-2d}$ où $d$ est la dimension relative de $p$.}
dans $\detaleplus X$ est le morphisme déduit
par adjonction du morphisme $\R p_!\Lambda(d)\cro{2d}\to\Lambda$
défini ci-dessus.
\end{definition2}

D'après \sga{4}{xviii}{3.2.4}, ce morphisme $\Cl p$ est un isomorphisme :
c'est la dualité de Poincaré.

\begin{proposition2}\label{16-proposition-composition-cl-lisses}
Si $f\colon Z\to Y$ et $g\colon Y\to X$ sont des morphismes lisses
compactifiables composables, le diagramme suivant
est commutatif dans $\detaleplus Z$ :
\[
\xymatrix{
\ar[rd]_{\Cl {j\circ i}}\Lambda\ar[r]^{\Cl i}  &
   i^?\Lambda \ar[d]^{i^?(\Cl j)} \\
&   i^?j^?\Lambda}
\]
\end{proposition2}

Ceci est énoncé en \sga{4}{xviii}{3.2.4} et résulte de la
compatibilité des morphismes traces à la composition,
cf.~propriété (Var~3) dans \sga{4}{xviii}{2.9}.

\begin{remarque2}
Si cette théorie avait été à notre disposition, il eût peut-être été plus
commode d'utiliser ici la construction des foncteurs $f^!$ pour $f$
lissifiable mentionnée dans l'introduction de \sga{4}{xviii}{0.4}.
Dans le cadre axiomatique des \guil{foncteurs homotopiques stables}, ceci
est réalisé dans \cite{THESE@Ayoub}.
\end{remarque2}

%>>>2
\subsection{Morphismes d'intersection complète lissifiables}%<<<2

\begin{definition2}\label{16-definition-intersection-complete}
Un morphisme d'intersection complète
est un morphisme 
$X\vers f S$ admettant localement une factorisation sous la forme
$X\vers i T\vers p S$ où $p$
est lisse et $i$ une immersion régulière (cf.~\sga{6}{vii}{1.4}).
On pose
$\dimrelvirt f = \dim p-\codim i$ : c'est la dimension relative virtuelle
de $f$ (cf.~\sga{6}{viii}{1.9}).
\end{definition2}

\begin{definition2}
On note $\Sch$ la catégorie dont les objets sont les $\ZZn$-schémas
quasi-compacts
admettant un faisceau inversible ample et dont les morphismes sont
les morphismes de type fini entre de tels schémas.
On note $\Schic$ la sous-catégorie de $\Sch$ ayant les mêmes objets mais
dont les morphismes sont les morphismes d'intersection complète.
\end{definition2}

Dans $\Sch$, tout morphisme $X\to Y$ peut se factoriser sous la forme
$X\vers i \PP^n_Y\vers \pi Y$ où $i$ est une immersion et $\pi$ la
projection canonique.
Tous les morphismes de $\Sch$ sont donc compactifiables, on
peut leur appliquer le formalisme des foncteurs $\R f_!$ et $f^!$.

Les morphismes de $\Schic$ admettent des factorisations globales dans
$\Schic$ sous la forme d'une immersion fermée régulière suivie d'un
morphisme lisse. 

\begin{definition2}
Pour tout morphisme $f\colon X\to Y$ dans $\Schic$, on peut
définir un foncteur
\[f^?\colon \detaleplus Y\to \detaleplus X\]
par la formule $f^?=f^!(-d)\cro{-2d}$
où $d=\dimrelvirt f$. 
\end{definition2}

Les foncteurs $f^?$ sont les foncteurs image inverse pour une structure de
catégorie fibrée convenable au-dessus de la catégorie $\Schic$ : on
utilisera implicitement les isomorphismes de transitivité $f^?g^?\simeq
(gf)^?$ associés à la composition de deux morphismes composables dans
$\Schic$.

\begin{definition2}\label{16-definition-cl-f}
Soit $f\colon X\to S$ un morphisme dans $\Schic$.
On suppose donnée une factorisation de $f$ dans $\Schic$ sous la forme
$X\vers i Y\vers p S$ où $i$ est une immersion régulière et $p$ un
morphisme lisse. On définit un morphisme
\[\Cl {p,i}\colon \Lambda\to f^?\Lambda\]
dans $\detaleplus X$ comme étant le morphisme composé
\[
\xymatrix{
\Lambda\ar[rd]_{\Cl {p,i}}\ar[r]^{\Cl i}&i^?\Lambda\ar[d]^{i^?(\Cl p)} \\
& f^?\Lambda}
\]
où $\Cl i$ est le morphisme de la définition~\ref{16-definition-cl-i}
et $\Cl p$ celui de la définition~\ref{16-definition-cl-p}.
\end{definition2}

\begin{theoreme2}\label{16-theoreme-independance-factorisation}
Soit $f\colon X\to S$ un morphisme dans $\Schic$.
Si $X\vers i Y\vers p S$ et $X\vers {i'}Y'\vers {p'}S$ sont
deux factorisations du type envisagé dans la
définition~\ref{16-definition-cl-f}, alors les deux morphismes suivants
dans la catégorie $\detaleplus X$ sont égaux :
\[\Cl {p,i}=\Cl{p',i'}\colon \Lambda\to f^?\Lambda\;\text{.}\]
\end{theoreme2}

La notation suivante s'avère assez commode pour cette démonstration :

\begin{definition2}
Si $f\colon Z\to Y$ et $g\colon Y\to X$ sont des morphismes composables
dans $\Schic$, $a\colon \Lambda\to g^?\Lambda$ et 
$b\colon \Lambda\to f^?\Lambda$
des morphismes dans $\detaleplus Y$ et $\detaleplus Z$ respectivement, on pose $a\bigstar
b=f^?(a)\circ b\colon \Lambda\to (g\circ f)^?\Lambda$.
\end{definition2}

Cette loi $\bigstar$ vérifiant une propriété d'associativité évidente, on
omettra les parenthèses.

Par définition, on a ainsi : $\Cl {p,i}=\Cl p\bigstar \Cl i$.
On veut vérifier l'égalité $\Cl p\bigstar\Cl i=\Cl {p'}\bigstar \Cl {i'}$.
Quitte à introduire le produit fibré de $Y$ et de $Y'$ au-dessus de $S$, on
peut supposer que \guil{$Y'$ coiffe $Y$}, à savoir qu'il existe un
morphisme lisse $q\colon Y'\to Y$ tel que $i=q\circ i'$ et $p'=p\circ q$ :
\[
\xymatrix{
& Y'\ar[rd]^{p'}\ar[d]^q & \\
X\ar[ru]^{i'}\ar[r]^i & Y\ar[r]^p &S}
\]
On a ainsi
\[\Cl{p',i'}=\Cl {p'}\bigstar \Cl{i'}=\Cl p\bigstar\Cl q\bigstar \Cl{i'}\;
\text{,}\]
la dernière égalité résultant de la
proposition~\ref{16-proposition-composition-cl-lisses}. On est ramené à montrer
l'égalité $\Cl i=\Cl q\bigstar\Cl {i'}$. Pour cela, on introduit le produit
fibré $X'$ de $X$ et $Y'$ au-dessus de $Y$ :
\[
\xymatrix{
\ar[d]^{q'}X'\ar[r]^j& Y'\ar[d]^q \\
X\ar@/^1pc/@{.>}[u]^s\ar[ru]_{i'}\ar[r]_i & Y}
\]
Le morphisme $i'$ donne naissance à la section $s$ de la projection
$q'\colon X'\to X$. Le morphisme $q$ étant lisse, l'immersion $j\colon
X'\to Y'$ est régulière. Admettons provisoirement les
égalités suivantes :
\[\Cl {q'}\bigstar \Cl s=\Id_\Lambda\;\text{,}\qquad 
\Cl q\bigstar \Cl j=\Cl i\bigstar \Cl {q'}\;\text{.}\]
Il vient :
\begin{eqnarray*}
\Cl i&=&\Cl i\bigstar \Cl {q'}\bigstar \Cl s\\
&=& \Cl q\bigstar \Cl j\bigstar \Cl s\;\text{.}
\end{eqnarray*}
On utilise alors la composition des morphismes de Gysin associés aux
immersions régulières 
(cf.~théorème~\ref{16-theoreme-composition-cl-immersions-regulieres}).
Celle-ci donne l'égalité $\Cl j\bigstar \Cl s=\Cl {i'}$ qui permet de
conclure que $\Cl i=\Cl q\bigstar \Cl {i'}$. Les deux lemmes qui suivent
permettent d'obtenir les deux égalités admises ci-dessus :

\begin{lemme2}\label{16-lemme-cb-lisse}
Soit un diagramme cartésien dans $\Sch$ :
\[
\xymatrix{
X'\ar[d]^{q'}\ar[r]^j & Y'\ar[d]^q \\
X\ar[r]^i & Y}
\]
On suppose que $q$ est lisse et que $i$ est une immersion régulière (donc
$j$ aussi). Alors on a l'égalité
\[\Cl q\bigstar \Cl j=\Cl i\bigstar\Cl {q'}\;\text{.}\]
\end{lemme2}

On peut supposer que $i$ est une immersion fermée.
On identifie $\Cl i$ (resp. $\Cl {j}$) à une classe dans 
$\HH^{2d}_X(Y,\Lambda(d))$ (resp. $\HH^{2c}_{X'}(Y',\Lambda(c))$) où $c$ est
la codimension de l'immersion régulière $i$. D'après la
proposition~\ref{16-proposition-classes-generalisees-cb}, on a
$q^\star(\Cl i)=\Cl j$. Bien que la vérification soit abracadabrante, 
la compatibilité du morphisme trace aux changements
de base (propriété (Var~2) de \sga{4}{xviii}{2.9}) permet de conclure.

\medskip

\begin{lemme2}
Soit $p\colon X\to S$ un morphisme lisse dans $\Sch$
admettant une section $s\colon
S\to X$ (qui est une immersion régulière). Alors, $\Cl p\bigstar \Cl
s=\Id_\Lambda$ dans $\detaleplus S$.
\end{lemme2}

Les endomorphismes de $\Lambda$ dans $\detaleplus S$ sont donnés par des sections
du faisceau $\Lambda$ dans $S$, il suffit de vérifier que les nombres
obtenus en passant aux points génériques de $S$ sont égaux à $1$. Comme on
peut supposer que $S$ est réduit et que la construction est compatible avec
le passage aux points génériques, on peut supposer que $S$ est le
spectre d'un corps $k$. Notons $x$ l'image de $\SP(k)$ dans $X$. Quitte à
remplacer $X$ par un voisinage ouvert, on peut supposer qu'il existe un
morphisme étale $\pi\colon X\to \Aff^d_k$ identifiant $x$ à l'image inverse
de l'origine dans $\Aff^d_k$. En utilisant l'isomorphisme évident
$\HH^{2d}_{(0,\dots,0)}(\Aff^d_k,\Lambda(d))\iso
\HH^{2d}_{x}(X,\Lambda(d))$, on se ramène au lemme suivant :

\begin{lemme2}
Pour tout entier naturel $d$ et tout schéma $S\in \Sch$, si on note $p\colon
\Aff^d_S\to S$ la projection et $s\colon S\to \Aff^d_S$ l'inclusion de
l'origine, on a l'égalité
\[\Cl p\bigstar \Cl s=\Id_\Lambda \]
dans $\detaleplus S$.
\end{lemme2}

L'énoncé est évident pour $d=0$. Une récurrence évidente s'appuyant sur le
théorème~\ref{16-theoreme-composition-cl-immersions-regulieres} et la
proposition~\ref{16-proposition-composition-cl-lisses} permet de se ramener au
cas où $d=1$, et comme précédemment, on peut supposer que $S=\SP(k)$ où
$k$ est un corps que l'on peut supposer séparablement clos. On se ramène
finalement au lemme suivant :

\begin{lemme2}\label{16-lemme-compatibilite-trace-c-un}
Pour tout corps séparablement clos $k$, si on note $p\colon \PP^1_k\to
\SP(k)$ la projection et $s\colon \SP(k)\to \PP^1_k$ l'inclusion de
$0$, on a l'égalité
\[\Cl p\bigstar \Cl s=\Id_\Lambda \]
dans $\detaleplus {\SP(k)}$.
\end{lemme2}

L'idéal de l'immersion fermée $s$ s'identifie au faisceau inversible
$\OO(-1)$. Par définition, l'image ${\Cl s}_{|\PP^1_k}$
de $\Cl s$ dans $\HH^2(\PP^1_k,\Lambda(1))$
est $c_1(\OO(-1))\cdot \Clf s$. Mais $\Clf s=-1$. Ainsi,
${\Cl s}_{|\PP^1_k}=c_1(\OO(1))$. Le degré du fibré en droites
$\OO(1)$ étant $1$, on peut conclure en utilisant la commutativité
du diagramme suivant (cf.~\sga{4}{xviii}{1.1.6}) :
\[
\xymatrix{
\ar[rd]_{\deg} \Pic(\PP^1_k) \ar[r]^-{c_1} &
\HH^2(\PP^1_k,\Lambda(1)) \ar[d]^{\Tr_p}_{\sim} \\
& \Lambda}
\]

\bigskip

\begin{definition2}\label{16-definition-cl-f-lissifiable}
Soit $f\colon X\to S$ un morphisme dans $\Schic$.
On note $\Cl f\colon
\Lambda\to f^?\Lambda$ le morphisme $\Cl {p,i}$ dans $\detaleplus X$ défini à
partir d'une factorisation de $f$ dans $\Schic$
sous la forme $f=p\circ i$ avec $i$ une
immersion régulière et $p$ un morphisme lisse. D'après le 
théorème~\ref{16-theoreme-independance-factorisation}, cette définition est
indépendante de la factorisation.
\end{definition2}

\begin{theoreme2}\label{16-theoreme-cl-f-transitivite}
Si $X\vers f Y$ et $Y\vers g Z$ sont des morphismes composables dans
$\Schic$, le diagramme suivant est commutatif dans $\detaleplus X$.
\[
\xymatrix{
\Lambda\ar[rd]_{\Cl {g\circ f}}\ar[r]^{\Cl f}&f^?\Lambda\ar[d]^{f^?(\Cl g)} \\
& (g\circ f)^?\Lambda }
\]
\end{theoreme2}

Paraphrasant \sga{6}{viii}{2.6}, on choisit une factorisation $Y\vers j
V'\vers p' Z$ dans $\Schic$ avec $j$ une immersion régulière et $p'$ lisse,
et une immersion régulière $X\vers i \PP^n_Y$, de façon à obtenir le
diagramme suivant :
\[
\xymatrix{
\ar[rd]_fX\ar[r]^i &\ar[d]^p \PP^n_Y\ar[r]^{j'} & \PP^n_{V'}\ar[d]^{p''}\\
& \ar[rd]_g Y\ar[r]^j & V' \ar[d]^{p'}\\
& & Z}
\]
En utilisant le
théorème~\ref{16-theoreme-composition-cl-immersions-regulieres} et la
proposition~\ref{16-proposition-composition-cl-lisses}, on obtient
\[\Cl{g\circ f} = (\Cl{p'}\bigstar\Cl{p''})\bigstar (\Cl {j'}\bigstar \Cl
i)
\;\text{.}\]
Le lemme~\ref{16-lemme-cb-lisse} donne l'égalité :
\[\Cl{p''}\bigstar \Cl {j'}=\Cl j\bigstar\Cl p\;\text{,}\]
ce qui permet d'obtenir :
\[\Cl {g\circ f}=(\Cl{p'}\bigstar \Cl j)\bigstar(\Cl p\bigstar \Cl i)
\;\text{,}\]
où l'on reconnaît l'égalité $\Cl {g\circ f}=\Cl g\bigstar \Cl f$.

\bigskip

\begin{proposition2}
Soit $f\colon X\to S$ un morphisme dans $\Schic$. On suppose que $f$ est
plat de dimension relative $d$. Alors le morphisme $\Cl f\colon \Lambda
\to f^?\Lambda$ correspond par adjonction au morphisme
$\R f_!\Lambda(d)[2d]\to \Lambda$ donné par le morphisme trace $\Tr_f\colon
\R^{2d}f_!\Lambda(d)\to \Lambda$.
\end{proposition2}

Compte tenu de la
proposition~\ref{16-proposition-compatibilite-sga-quatre-et-demi}, cela
résulte de \sgaquatreetdemi{Cycle}{2.3.8~(i)}.

\begin{remarque2}
Si $f\colon X\to Y$ est un morphisme propre
dans $\Schic$ de dimension relative
virtuelle $d$, le morphisme $\Cl f$ permet de définir, pour tout $K\in
\detaleplus Y$, un morphisme $f_\star\colon \HH^{p}(X,f^\star K)\to
\HH^{p-2d}(Y,K(-d))$, compatible à la composition. On peut aussi en définir une
version à supports $f_\star\colon \HH^{p}_Z(X,f^\star K)\to
\HH^{p-2d}_{Z'}(Y,K(-d))$ dès que $Z$ et $Z'$ sont des fermés de
$X$ et $Y$ respectivement tels que $f(Z)\subset Z'$.
\end{remarque2}

%>>>2
%>>>1
\section{Théorème de pureté}%<<<1
\label{16-section-purete}
\subsection{Énoncés}%<<<2
L'objectif de cette section est de donner une démonstration
du théorème suivant :

\begin{theoreme2}\label{16-theoreme-purete-absolue}
Soit $X$ un $\ZZn$-schéma régulier. Soit $Y$ un sous-schéma (fermé)
de $X$ qui est aussi
régulier. On note $i\colon  Y\to X$ l'immersion, et $c$ sa codimension.
Alors, le morphisme de
Gysin $\Cl i\colon
\Lambda\to i^?\Lambda=i^!\Lambda(c)\cro{2c}$ est un isomorphisme dans
$\detaleplus Y$.
\end{theoreme2}

\begin{corollaire2}\label{16-corollaire-purete-absolue}
Soit $f\colon X\to S$ un morphisme de type fini entre $\ZZn$-schémas
réguliers. On suppose que $X$ et $S$ admettent un faisceau ample. Alors, le
morphisme de Gysin $\Cl f\colon \Lambda(d)[2d]\to f^!\Lambda$ est un
isomorphisme dans $\detaleplus X$, où $d$ désigne la dimension relative virtuelle
de $f$.
\end{corollaire2}

\begin{corollaire2}\label{16-corollaire-purete-absolue-dcn}
Soit $X$ un $\ZZn$-schéma régulier. Soit $D$ un diviseur à croisements
normaux dans $X$. On note $j\colon X-D\to X$ l'inclusion de son 
complémentaire. Alors, $\R j_\star \Lambda$ appartient à $\detalebctf X$.
Plus précisément, si $D=D_1+\dots+D_n$ est un diviseur à croisements
normaux strict, alors $\R^1 j_\star \Lambda$ s'identifie à $\bigoplus_{1\leq
i\leq n} \Lambda_{D_i}(-1)$ et $\R^\star j_\star \Lambda$ est l'algèbre
extérieure sur $R^1 j_\star \Lambda$.
\end{corollaire2}

Ce corollaire mérite une démonstration. Pour la première assertion, on peut
travailler localement pour la topologie étale sur $X$ ; il suffit donc
d'établir la deuxième assertion. On suppose que $D=D_1+\dots+D_n$ est un
diviseur à croisements normaux strict. On note $j_i\colon X-D_i\to X$
l'inclusion du complémentaire de $D_i$ pour tout $i$. Nous allons montrer
que le morphisme de Künneth
\[\R j_{1\star}\Lambda \otimes^{\L}\dots \otimes^{\L} \R
j_{n\star}\Lambda\to \R j_\star\Lambda\]
est un isomorphisme dans $\detale X$, ce qui impliquera le résultat vu que
les faisceaux $\R^q j_{i\star}\Lambda$ sont connus par pureté ($\Lambda$
pour $q=0$, $\Lambda_{D_i}(-1)$ pour $q=1$ et $0$ sinon) et qu'ils sont
plats.

On procède par récurrence sur $n$. Les cas $n=0$ et $n=1$ sont évidents. On
suppose $n\geq 2$, on pose $D'=D_2+\dots+D_n$ et on fait l'hypothèse que le
résultat est connu pour $D'$. Il s'agit donc de monter que si on note
$j'\colon X-D'\to X$ l'inclusion du complémentaire de $D'$, alors le
morphisme de Künneth
\[\R j_{1\star}\Lambda\otimes^{\L} \R j'_\star \Lambda\to \R j_\star\Lambda\]
est un isomorphisme. Autrement dit, le morphisme canonique
\[\R\SheafHom(\Lambda_{X-D_1},\Lambda)\otimes^{\L}
\R\SheafHom(\Lambda_{X-D'},\Lambda)\to \R\SheafHom(\Lambda_{X-D_1}\otimes
\Lambda_{X-D'},\Lambda)\]
est un isomorphisme dans $\detale X$. À $K$ (resp. $L$) fixé dans $\detale
X$, la famille des $L$ (resp. $K$) tels que le morphisme
\[\R\SheafHom(K,\Lambda)\otimes^{\L}
\R\SheafHom(L,\Lambda)\to \R\SheafHom(K\otimes^{\L} L,\Lambda)\]
soit un isomorphisme, propriété que nous appellerons (Kü), est une
sous-catégorie triangulée de $\detale X$.

Pour $K=\Lambda$ ou $L=\Lambda$, la condition (Kü) est évidemment vérifiée,
la montrer pour $(\Lambda_{X-D_1},\Lambda_{X-D'})$ revient donc, par
dévissage, à la montrer pour $(\Lambda_{D_1},\Lambda_{X-D'})$ ou encore
pour $(\Lambda_{D_1},\Lambda_{D'})$. Il résulte aussitôt du théorème de
pureté et des compatibilités obtenues que si $Y$ et $Z$ sont deux
sous-schémas fermés réguliers de $X$ s'intersectant transversalement
(\emph{i.e.} $Y\cap Z$ est régulier de codimension la somme des
codimensions de $Y$ et de $Z$), alors $(\Lambda_Y,\Lambda_Z)$ vérifie (Kü).
En particulier, $(\Lambda_{D_1},\Lambda_{D_i})$ vérifie (Kü) pour $i\geq 2$
et plus généralement, pour tout sous-ensemble non vide $I$ de
$\{2,\dots,n\}$, $(\Lambda_{D_1},\Lambda_{D_I})$ vérifie (Kü) où $D_I$ est
l'intersection des $D_i$ pour $i\in I$. En utilisant la suite exacte
standard
\[0\to \Lambda_{D'}\to \bigoplus_{2\leq i\leq n}\Lambda_{D_i}\to
\bigoplus_{2\leq i<j\leq n} \Lambda_{D_{ij}}\to \dots\;\text{,}\]
on en déduit par dévissage la condition (Kü) pour
$(\Lambda_{D_1},\Lambda_{D'})$, ce qu'il fallait démontrer.

\begin{definition2}
Un couple régulier est un couple $(X,Y)$ où $X$ est un $\ZZn$-schéma régulier
et $Y$ un sous-schéma fermé de $X$ qui est régulier. On dit que $(X,Y)$ est
pur si la conclusion du théorème~\ref{16-theoreme-purete-absolue} est vraie
pour l'inclusion de $Y$ dans $X$. Si $\overline y\to Y$ est un point
géométrique de $Y$, on dira que $(X,Y)$ est pur en $\overline y$ si
$(\Cl i)_{\overline y}$ est un isomorphisme dans $\detaleplus {\overline y}$.
\end{definition2}

Le théorème~\ref{16-theoreme-purete-absolue} peut ainsi se reformuler en
disant que tout couple régulier est pur. Dans la
sous-section~\ref{16-subsection-purete-ponctuelle} sera introduite la notion de
pureté ponctuelle qui consiste à étudier les couples réguliers de la forme
$(X,x)$ où $X$ est un schéma local régulier de
point fermé $x$. Pour
démontrer le théorème de pureté, il suffira de savoir que les couples
réguliers de cette forme sont purs. Dans la
sous-section~\ref{16-subsection-changement-de-coefficients}, on se ramènera
au cas où l'anneau de coefficients $\Lambda$ est $\ZZ/\ell\ZZ$ avec $\ell$
un nombre premier inversible sur les schémas réguliers considérés. Dans
la sous-section~\ref{16-subsection-diviseurs-reguliers}, on établira quelques
propriétés utiles concernant la pureté des couples réguliers donnés par des
diviseurs. 
Comme dans la démonstration de \cite{AZUMINO@Fujiwara},
la démonstration de la pureté ponctuelle pour des schémas réguliers
arbitraires se ramènera à celle des
schémas réguliers qui sont de type fini sur
un trait $S$ (d'inégale caractéristique). Dans la
sous-section~\ref{16-subsection-schemas-sur-un-trait}, on obtiendra des
conditions suffisantes pour montrer que des schémas réguliers de type fini
sur $S$ sont ponctuellement purs. La
sous-section~\ref{16-subsection-geometrie-logarithmique} donnera les énoncés
de géométrie logarithmique permettant d'établir que si $(X,M)$ est un
log-schéma log-lisse sur un trait (muni de sa log-structure canonique)
et que le schéma $X$ est régulier, alors $X$ est ponctuellement pur.
La démonstration du théorème~\ref{16-theoreme-purete-absolue}
sera donnée dans la
sous-section~\ref{16-subsection-demonstration-theoreme-purete}. Elle utilisera
les résultats des sous-sections précédentes ainsi que trois théorèmes de
résolution des singularités que l'on peut résumer ainsi :
\begin{itemize}
\item utilisation d'altérations pour obtenir un schéma à réduction
semi-stable à partir d'un schéma (normal) sur $S$
(cf.~\cite[proposition~4.4.1]{BS@Vidal}) ;
\item résolution des singularités d'une action modérée
d'un groupe fini sur un log-schéma log-régulier de façon à
obtenir une action très modérée (cf. exposé VI) ;
\item résolution des log-singularités des log-schémas log-réguliers
(cf.~exposé VI).
\end{itemize}

%>>>2
\subsection{Pureté ponctuelle}%<<<2
\label{16-subsection-purete-ponctuelle}

\begin{definition2}
Soit $X$ un $\ZZn$-schéma local régulier.
On dit que $X$ est ponctuellement pur en
son point fermé $x$ si le morphisme $\Cl i\colon \Lambda\to i^?\Lambda$
est un isomorphisme dans $\detaleplus x$ où $i\colon x\to X$ est
l'inclusion du point fermé de $X$.
\end{definition2}

Un schéma local régulier est ponctuellement pur en son point fermé si et
seulement si son hensélisé (resp. son hensélisé strict) l'est.

\begin{definition2}
Soit $X$ un $\ZZn$-schéma. Si $x\in X$, on dit que $X$ est ponctuellement
pur au point $x$ si le localisé de $X$ en $x$ est ponctuellement pur en son
point fermé. On dit que $X$ est ponctuellement pur s'il l'est en tous ses
points.
\end{definition2}

La proposition suivante est \cite[proposition~2.2.4]{AZUMINO@Fujiwara}. La
démonstration de cet article semble compliquée puisqu'elle passe par des
résultats plus fins que ceux dont nous avons besoin. On en redonne donc
une démonstration plus courte.

\begin{proposition2}\label{16-proposition-deux-sur-trois}
Soit $i\colon Y\to X$ une immersion fermée entre schémas réguliers. Le
nombre de conditions satisfaites parmi les trois suivantes ne peut pas être
deux :
\begin{enumerate}
\item[(a)] Le couple régulier $(X,Y)$ est pur ;
\item[(b)] Le schéma $Y$ est ponctuellement pur ;
\item[(c)] Le schéma $X$
est ponctuellement pur aux points situés dans l'image de $Y$.
\end{enumerate}
\end{proposition2}

Soit $y\in Y$, notons $V(y)$ le localisé de $Y$ en $y$ et $V(x)$ celui de
l'image $x$ de $y$ dans $X$. On a un diagramme de schémas :
\[
\xymatrix{
y\ar[rd]_{i_x}\ar[r]^{i_y} & V(y)\ar[d]^{i'} \\
& V(x)}
\]
La composition des morphismes de Gysin donne le diagramme commutatif
suivant dans $\detaleplus y$ :
\[
\xymatrix{
\ar[rrd]_{\Cl {i_x}}\Lambda\ar[rr]^{\Cl {i_y}}
&   & i_y^?\Lambda\ar[d]^{i^?_y\Cl {i'}} \\
&  & i_x^?\Lambda}
\]

Sur ce diagramme, on voit aussitôt que (a) et (b) impliquent (c) et que
(a) et (c) impliquent (b). Montrons que (b) et (c) impliquent (a). Il
s'agit de montrer que pour tout point $y$ de $Y$, le morphisme 
$i_y^\star \Cl {i'}$ est un isomorphisme.
On peut procéder par récurrence sur la
dimension de $V(y)$. On peut ainsi supposer que le support d'un
cône $C$ du morphisme
$\Cl {i'}$ dans $\detaleplus {V(y)}$ est contenu dans $\left\{y\right\}$.
Mézalor, le morphisme
canonique $i_y^!C\to i_y^\star C$ est un isomorphisme ; le diagramme
ci-dessus montre que $i_y^!C=0$, ce qui permet de conclure que $C=0$ et
finalement d'obtenir (a).

\bigskip

Rappelons quelques propriétés importantes concernant la pureté ponctuelle :

\begin{proposition2}[{\cite[proposition~2.2.2]{AZUMINO@Fujiwara}}]%
\label{16-proposition-purete-ponctuelle-completion}
Soit $X$ un schéma local strictement hensélien régulier.
Le complété $\hat{X}$ est
ponctuellement pur en son point fermé si et seulement si $X$ l'est.
\end{proposition2}

\begin{proposition2}[{\cite[corollary~2.2.3]{AZUMINO@Fujiwara}}]%
\label{16-proposition-purete-egale-caracteristique}
Soit $k$ un corps premier. Si $X$ est schéma régulier qui est
un $k$-schéma, alors $X$ est ponctuellement pur.
\end{proposition2}

%>>>2
\subsection{Changement de coefficients}%<<<2
\label{16-subsection-changement-de-coefficients}

\begin{proposition2}\label{16-proposition-lemme-chinois}
Soit $n$ un entier naturel non nul.
Soit $n=\prod_{j=1}^k \ell_j^{\nu_j}$ la factorisation de $n$ en produit de
puissances de nombres premiers distincts. Un couple régulier $(X,Y)$ est
pur relativement à l'anneau de coefficients $\ZZ/n\ZZ$ si et seulement
s'il l'est relativement à l'anneau de coefficients $\ZZ/\ell_j^{\nu_j}\ZZ$
pour tout $j\in\{1,\dots,k\}$.
\end{proposition2}

Cela résulte aussitôt du lemme chinois et du fait que si $m$ est un
entier naturel divisant $n$, alors pour toute immersion fermée régulière
$i\colon Y \to X$, le diagramme évident commute dans $\detalepluscoefs{Y}{\ZZ/n\ZZ}$ :
\[
\xymatrix{
\ZZ/n\ZZ\ar[r]^{\Cl i} \ar[d]& i^?\ZZ/n\ZZ\ar[d]\\
\ZZ/m\ZZ\ar[r]^{\Cl i}       & i^?\ZZ/m\ZZ}
\]

\bigskip

\begin{proposition2}\label{16-proposition-ops-n-egale-ell}
Soit $\ell$ un nombre premier. Pour tout entier $\nu\geq 1$, un couple
régulier $(X,Y)$ est pur relativement à l'anneau de coefficients
$\ZZ/\ell\ZZ$ si et seulement s'il l'est relativement à l'anneau de
coefficients $\ZZ/\ell^\nu\ZZ$.
\end{proposition2}

En utilisant la résolution de Godement des faisceaux $\ZZ/\ell^\nu\ZZ(c)$
(où $c$ est la codimension de l'immersion $i\colon Y\to X$) pour tout
$\nu$, on peut représenter les morphismes de Gysin
$\Cl i\colon \ZZ/\ell^\nu\ZZ\to i^?\ZZ/\ell^\nu\ZZ$ dans
$\detalepluscoefs{Y}{\ZZ/\ell^\nu\ZZ}$ par des cocycles. Un tel cocycle pour
$\nu_0$ fixé induit pour tout entier $\nu\leq \nu_0$ un cocycle
représentant le morphisme de Gysin à coefficients dans $\ZZ/\ell^\nu\ZZ$.
Les propriétés élémentaires de la résolution de Godement font que, si on le
souhaite, on peut en fait trouver une famille compatible de cocycles pour
tout $\nu\in\NN$.

Compte tenu de ces observations, une fois ces cocycles convenablement
choisis, on dispose d'un cône privilégié $C(\nu)$ du morphisme $\Cl
i\colon\ZZ/\ell^\nu\ZZ\to i^?\ZZ/\ell^\nu\ZZ$ dans
$\detalepluscoefs{Y}{\ZZ/\ell^\nu\ZZ}$ pour tout $\nu\in\NN$ et de triangles
\[
\xymatrix{
C(\mu)\ar[r] & C(\mu+\nu)\ar[r] & C(\nu) \ar[r] & C(\mu)[1] }
\]
dans
$\detalepluscoefs{Y}{\ZZ/\ell^{\mu+\nu}\ZZ}$ pour tous $(\mu,\nu)\in\NN^2$.

Par conséquent, si $C(1)=0$, il vient que pour tout $\nu\geq 1$,
$C(\nu)=0$. Inversement, si $C(1)$ est non nul, son premier objet de
cohomologie non nul s'injecte dans celui de $C(\nu)$ pour tout $\nu\geq 1$.

%>>>2
\subsection{Diviseurs réguliers}%<<<2
\label{16-subsection-diviseurs-reguliers}

\begin{definition2}
Si $X$ est un schéma et $\overline x \to X$ un point géométrique, on note
$V(\overline x)$ l'hensélisé strict de $X$ en $\overline x$ et 
$i_{\overline x}\colon V(\overline x)\to X$ le morphisme canonique.
\end{definition2}

\begin{proposition2}\label{16-proposition-caracterisation-purete-diviseur}
Soit $X$ un schéma régulier. Soit $D$ un diviseur régulier de $X$. Le
couple régulier $(X,D)$ est pur si et seulement si pour tout point
géométrique $\overline x\to D$, on a
$\HH^q_\et(V(\overline x)-i_{\overline x}^{-1}(D),\Lambda)=0$ pour tout
$q\geq 2$.
\end{proposition2}

Cela résulte du calcul de
$\HH^q_{\et}(V(\overline x)-i_{\overline x}^{-1}(D),\Lambda)$ pour
$q\in\{0,1\}$ (cf.~\sgaquatreetdemi{Cycle}{2.1.4}).

\bigskip

\begin{proposition2}\label{16-proposition-purete-argument-de-transfert-facile}
On suppose que l'anneau de coefficients est $\ZZ/\ell\ZZ$ où $\ell$ est un
nombre premier.
Soit $f\colon Y\to X$ un morphisme fini et plat de degré constant premier à
$\ell$ entre $\ZZn[\ell]$-schémas réguliers. Soit $D$ un diviseur régulier
de $X$. On suppose que $D'=f^{-1}(D)_{\red}$ est un diviseur régulier de
$Y$. Si le couple régulier $(Y,D')$ est pur, alors $(X,D)$ aussi.
\end{proposition2}

Grâce à la proposition~\ref{16-proposition-caracterisation-purete-diviseur},
on peut choisir un point géométrique de $D$ et remplacer $X$ par son
hensélisé strict en ce point. On suppose donc que $X$ et $D$ sont locaux
strictement henséliens et on se concentre sur la pureté du couple $(X,D)$
en le point fermé de $D$.
Le schéma $Y$ est
alors réunion disjointe finie de schémas locaux strictement henséliens ; au
moins un de ceux-ci est de degré premier à $\ell$ sur $X$. On peut donc
supposer que $Y$ aussi est local strictement hensélien.
Il suffit alors de montrer que $\HH^q(X-D,\ZZ/\ell\ZZ)$ s'injecte dans
$\HH^q(Y-D',\ZZ/\ell\ZZ)$, ce qui résulte du lemme suivant :

\begin{lemme2}
On suppose que l'anneau de coefficients $\Lambda$
est $\ZZ/\ell\ZZ$ où $\ell$ est un
nombre premier.
Soit $f\colon Y\to X$ un morphisme de présentation finie, fini et plat de
rang $d$ premier à $\ell$ entre $\ZZn[\ell]$-schémas.
Alors,
le morphisme canonique $\Lambda\to f_\star \Lambda$ est un monomorphisme
scindé dans $\detaleplus X$.
\end{lemme2}

D'après \sga{4}{xvii}{6.2.3}, on a un morphisme $\Tr_f\colon
f_\star\Lambda\to\Lambda$ tel que la composée
\[\Lambda\to f_\star\Lambda\to\Lambda\]
soit la multiplication par $d$, ce qui donne le scindage voulu puisque $d$
est inversible dans $\Lambda$.

\bigskip

\begin{proposition2}\label{16-proposition-up-and-down}
On suppose que l'anneau de coefficients $\Lambda$
est $\ZZ/\ell\ZZ$ où $\ell$ est un
nombre premier. Soit $X$ un $\ZZn[\ell]$-schéma régulier.
Soit $f$ une fonction sur $X$
dont le lieu des zéros $D=V(f)$ soit un diviseur régulier de $X$. On pose
$X'=\SP\left(\OO_X[T]/(T^\ell-f)\right)$. On note $\pi\colon X'\to X$ la projection,
$D'=\pi^{-1}(D)_{\red}$ (noter que $D'\to D$ est un isomorphisme). Alors,
$X'$ est un schéma régulier, et le couple régulier $(X',D')$ est pur si et
seulement si le couple régulier $(X,D)$ l'est.
\end{proposition2}

Soit $\overline x$ un point géométrique de $D$ (on identifiera aussi
$\overline x$ à un point géométrique de $D'$).
On va en fait montrer que
$(X,D)$ est pur en $\overline x$ si et seulement si $(X',D')$
l'est. On peut supposer que $X$ est le spectre premier d'un anneau local
strictement hensélien $A$ d'idéal maximal $\got{m}$ et que $\overline
x$ est au-dessus du point fermé de $X$. On a évidemment $f\in\got{m}$ ;
le fait que $D=V(f)$ soit régulier revient à dire que $f\not\in 
{\got{m}}^2$.

Notons $A'=A[T]/(T^\ell-f)$. En considérant le déterminant de
l'endomorphisme de $A'$ comme $A$-module donné par la multiplication par un
élément $b\in A'$, on observe que $b$ est inversible dans $A'$ si et
seulement si son image dans l'algèbre locale $(A/\got{m})[T]/(T^\ell)$
est inversible. Il en résulte que $A'$ est local d'idéal maximal
$\got{m}'=(T)+\got{m}A'$. Par ailleurs, on a un isomorphisme $A/(f)\iso
A'/(T)$ (\emph{i.e.} $D'\to D$ est un isomorphisme). On construit facilement un
isomorphisme $A_{(f)}[T]/(T^d-f)\iso A'_{(T)}$, ce qui montre que le
localisé de $A'$ en $(T)$ est un anneau de valuation discrète. La
codimension de l'idéal premier $(T)$ dans $A'$ est donc $1$. Compte tenu du
fait que $A'/(T)$ soit régulier, il en résulte que $X'$ est régulier.

On peut considérer, pour tout entier $\nu\geq 0$, le
$X$-schéma affine $X^\nu=\SP\left(\OO_X[T]/(T^{\ell^\nu}-f)\right)$. En élevant $T$ à
la puissance $\ell$, on obtient une tour de morphismes
\[\dots \to X^{\nu+1}\to X^{\nu} \to \dots \to X^1\to X^0\;\text{,}\]
le dernier morphisme $X^1\to X^0$ s'identifiant à $\pi\colon X'\to X$.
Au-dessus de $X-D$, cette tour de morphismes définit une tour de revêtements
étales galoisiens de $X-D$ de groupe de Galois
$\ZZ_\ell(1)=\underset{\nu}\lim\,\mu_{\ell^\nu}$ où pour tout entier $n$
inversible dans $A$, on note simplement $\mu_n=\mu_n(A)$.

Cette tour de revêtements définit un morphisme (surjectif) de groupes
profinis $\pi_1^\et(X-D)^\ab\to \ZZ_\ell(1)$. Par conséquent, on a un morphisme
de topos $\rho_f\colon (X-D)_\et\to \BB\ZZ_\ell(1)$ où $\BB\ZZ_\ell(1)$
désigne le topos des $\ZZ_\ell(1)$-ensembles discrets.

\begin{lemme2}\label{16-lemme-interpretation-purete-diviseur}
Le couple régulier $(X,D)$ est pur en $\overline x$
si et seulement si le morphisme
\[\RG(\BB\ZZ_\ell(1),\mu_\ell)\to\RG((X-D)_\et,
\mu_\ell)\]
induit par le morphisme de topos $\rho_f$ est un isomorphisme dans la
catégorie dérivée des groupes abéliens.
\end{lemme2}

Ce lemme découle des deux lemmes suivants :

\begin{lemme2}
Pour tout entier $q\geq 2$, $\HH^q(\BB\ZZ_\ell(1),\mu_\ell)=0$ et on a
des isomorphismes canoniques
\[\HH^0(\BB\ZZ_\ell(1),\mu_\ell)\simeq \mu_\ell\;\text{,}\quad
\HH^1(\BB\ZZ_\ell(1),\mu_\ell)\simeq  \Hom(\ZZ_\ell(1),\mu_\ell)
\simeq \ZZ/\ell\ZZ\;\text{.}\]
\end{lemme2}

Il s'agit de montrer que $\ZZ_\ell(1)$ est de $\ell$-dimension
cohomologique $1$. Pour cela, voir par exemple
\cite[§3.4, chapitre~I]{CG@Serre}.

\begin{lemme2}
Le morphisme composé \[\ZZ/\ell\ZZ\simeq
\HH^1(\BB\ZZ_\ell(1),\mu_\ell)\vers{\rho_f^\star}\HH^1((X-D)_\et,\mu_\ell)\iso
\HH^2_D(X,\mu_\ell)\] est donné (au signe près) par la classe de Gysin
$\Cl{D\subset X}$.
\end{lemme2}

Un relèvement dans $\HH^1(X-D,\mu_\ell)$ de $\Cl{D\subset X}$ est donné par
l'image de $f$ par le morphisme de bord $\HH^0((X-D)_\et,\Gm)\to
\HH^1((X-D)_\et,\mu_\ell)$ associé à la suite exacte de Kummer
\[0\to \mu_\ell\to \Gm\vers{p} \Gm\to 0\;\text{,}\]
autrement dit par le $\mu_\ell$-torseur $p^{-1}(f)\subset \Gm$.
Géométriquement, ce
torseur s'identifie tautologiquement au revêtement galoisien $X'-D'\to X-D$
de groupe de Galois $\mu_\ell$. Bien sûr,
la classe de ce $\mu_\ell$-torseur est 
donnée par le morphisme évident $\pi_1^\et(X-D)^\ab\to\ZZ_\ell(1)\to
\mu_\ell$ donnant l'image par $\rho_f^\star$ du générateur canonique de
$\HH^1(\BB\ZZ_\ell(1),\mu_\ell)$.

\bigskip

On peut appliquer le lemme~\ref{16-lemme-interpretation-purete-diviseur} à
$X'$ : il vient que le couple régulier $(X',D')$ est pur en $\overline x$
si et seulement si
le morphisme
\[\RG(\BB\ZZ_\ell(1),\mu_\ell)\to\RG((X'-D')_\et,
\mu_\ell)\]
induit par le morphisme de topos $\rho_T\colon (X'-D')_\et\to
\BB\ZZ_\ell(1)$ est un isomorphisme.

On dispose d'un carré commutatif de topos :
\[
\xymatrix{
\ar[d]^g (X'-D')_\et\ar[r]^{\rho_T} & \BB\ZZ_\ell(1)
 \ar[d]_{g'}\\
(X-D)_\et\ar[r]^{\rho_f} & \BB\ZZ_\ell(1)}
\]
où $g$ est induit par $\pi\colon X'\to X$ et $g'$ par la multiplication par
$\ell$ sur $\ZZ_\ell(1)$. Le faisceau $g_\star \ZZ/\ell\ZZ$
s'identifie canoniquement à
$\rho_f^\star g'_\star(\ZZ/\ell\ZZ)$. Il en découle aisément que le couple
régulier $(X',D')$ est pur en $\overline x$
si et seulement si le morphisme canonique
\[\RG(\BB\ZZ_\ell(1),g'_\star(\ZZ/\ell\ZZ))\to
\RG((X-D)_\et,\rho_f^\star
g'_\star(\ZZ/\ell\ZZ)) \]
est un isomorphisme.

La cohomologie relative du morphisme de topos $\rho_f\colon (X-D)_\et\to
\BB\ZZ_\ell(1)$ définit un foncteur triangulé 
\[F\colon \dplus(\BB\ZZ_\ell(1),\ZZ/\ell\ZZ)\to
\dplus(\ZZ/\ell\ZZ)\]
tel que pour tout $K\in\dplus(\BB\ZZ_\ell(1),\ZZ/\ell\ZZ)$, $F(K)$ soit
isomorphe à un cône du morphisme canonique $\RG(\BB\ZZ_\ell(1),K)\to
\RG((X-D)_\et,\rho_f^\star K)$.

Le lemme suivant découle de ce qui précède :

\begin{lemme2}
Le couple régulier $(X,D)$ est pur en $\overline x$
si et seulement si $F(\ZZ/\ell\ZZ)=0$
tandis que $(X',D')$ est pur en $\overline x$
si et seulement si $F(g'_\star(\ZZ/\ell\ZZ))=0$.
\end{lemme2}

Comme $\ZZ_\ell(1)$ est un pro-$\ell$-groupe, le faisceau
$g'_\star(\ZZ/\ell\ZZ)$ est une extension successive de $l$ copies de
$\ZZ/\ell\ZZ$. Le foncteur $F$ étant triangulé, on en déduit aussitôt que
si $F(\ZZ/\ell\ZZ)$ est nul, alors $F(g'_\star\ZZ/\ell\ZZ)$ aussi et que si
$F(\ZZ/\ell\ZZ)$ est non nul, son premier objet de cohomologie non nul
s'injecte dans celui de $F(g'_\star(\ZZ/\ell\ZZ))$. Ceci achève la
démonstration de la proposition~\ref{16-proposition-up-and-down}.

%>>>2
\subsection{Schémas sur un trait}%<<<2
\label{16-subsection-schemas-sur-un-trait}

Soit $S$ un trait. On note $s$ son point fermé, $\eta$ son point
générique et $\pi$ une uniformisante.

\begin{proposition2}\label{16-proposition-un-trait-est-ponctuellement-pur}
Le trait $S$ est ponctuellement pur.
\end{proposition2}

Il s'agit de montrer que $S$ est ponctuellement pur en son point fermé.
On peut supposer que $S$ est strictement hensélien ; cela résulte alors
facilement du fait que le corps des fractions de $S$ soit de $\ell$-dimension
cohomologique $1$ pour tout nombre premier inversible sur $S$
(cf.~\sga{4}{x}{2.2}).

\begin{proposition2}
Pour tout entier naturel $n$, l'espace affine $\Aff^n_S$ est ponctuellement
pur.
\end{proposition2}

D'après la proposition~\ref{16-proposition-purete-egale-caracteristique}, les
schémas $\Aff^n_s$ et $\Aff^n_\eta$ sont ponctuellement purs. Ainsi,
$\Aff^n_S$ est ponctuellement pur en les points de la fibre générique.
Pour établir la pureté ponctuelle de $\Aff^n_S$ en les
points de la fibre spéciale, on utilise la
proposition~\ref{16-proposition-deux-sur-trois} : il suffit de montrer que
le couple régulier $(\Aff^n_S,\Aff^n_s)$ est pur.
Le cas $n=0$ résulte de la
proposition~\ref{16-proposition-un-trait-est-ponctuellement-pur} et le cas
général en découle en vertu du théorème de changement de base lisse.

\begin{corollaire2}\label{16-corollaire-lisse-implique-ponctuellement-pur}
Un $S$-schéma lisse est ponctuellement pur.
\end{corollaire2}

\begin{definition2}
Soit $p\colon X\to S$ un morphisme de type fini, avec $X$
régulier et admettant un faisceau ample.
On pose $K_X=p^?\Lambda_S$ et on dispose d'un morphisme de
Gysin $\Cl{X/S} \colon \Lambda_X\to K_X$ dans $\detaleplus X$
(cf. définition~\ref{16-definition-cl-f-lissifiable}).
\end{definition2}

\begin{proposition2}
Soit $p\colon X\to S$ un morphisme de type fini, avec $X$
régulier et admettant un faisceau ample.
Le schéma $X$ est ponctuellement pur si et seulement si le morphisme
$\Cl{X/S}\colon \Lambda_X\to K_X$ est un isomorphisme dans $\detaleplus X$.
\end{proposition2}

On choisit une factorisation $X\vers i Y \vers q S$ de $p$ (dans la
catégorie $\Schic$) avec $Y$ lisse
sur $S$ et $i$ une immersion fermée (régulière). D'après le
théorème~\ref{16-theoreme-cl-f-transitivite} (ou plutôt par définition de
$\Cl{X/S}$), le diagramme suivant est
commutatif :
\[
\xymatrix{
\ar[rd]_{\Cl{X/S}}\Lambda\ar[r]^{\Cl i} & i^?\Lambda
\ar[d]^{i^?\Cl q}_\sim\\
& i^?q^?\Lambda}
\]
Le morphisme $q$ étant lisse, le morphisme de Gysin $\Cl q$ est un
isomorphisme. Par conséquent, $\Cl {X/S}$ est un isomorphisme si et
seulement si $\Cl i\colon \Lambda\to i^?\Lambda$ en est un. D'après la
proposition~\ref{16-proposition-deux-sur-trois} et compte tenu du fait que $Y$
soit ponctuellement pur (cf. corollaire~%
\ref{16-corollaire-lisse-implique-ponctuellement-pur}), ceci équivaut encore
à dire que $X$ est ponctuellement pur.

\begin{corollaire2}\label{16-corollaire-lisse-sur-ponctuellement-pur}
Soit $X$ un $S$-schéma de type fini qui est régulier. Soit $Y$ un
$X$-schéma lisse. Si $X$ est ponctuellement pur, alors $Y$ aussi.
\end{corollaire2}

\begin{proposition2}\label{16-proposition-purete-argument-de-transfert}
Soit $f\colon X\to Y$ un morphisme propre et dominant de $S$-schémas où $X$
et $Y$ sont supposés de type fini sur $S$, intègres, réguliers et
admettant des faisceaux amples.
On suppose de plus que $f$ est génériquement étale de degré $d$ inversible
dans $\Lambda$. Alors, la pureté ponctuelle de $X$ implique celle de $Y$.
\end{proposition2}

Le morphisme $f$ est localement d'intersection complète lissifiable de
dimension relative virtuelle zéro, d'où $f^?=f^!$. Le morphisme de Gysin
relatif à $f$ est un morphisme $\Cl f\colon \Lambda\to f^!\Lambda$.

\begin{lemme2}
On peut généraliser le morphisme $\Cl f\colon \Lambda\to f^!\Lambda$ en des
morphismes $f^\star M\to f^!M$, fonctoriellement en $M\in\detaleplus Y$.
\end{lemme2}

Le morphisme $\Cl f\colon \Lambda\to f^!\Lambda$ correspond par adjonction
à un morphisme $\R f_!\Lambda\to \Lambda$,
que l'on peut tensoriser avec $M$
pour obtenir (\emph{via} la formule de projection) un morphisme
$\R f_!f^\star M\to M$
qui correspond lui-même par adjonction au morphisme
$f^\star M\to f^!M$ du type recherché.

\medskip

En appliquant la fonctorialité de la construction du lemme
au morphisme $\Cl{Y/S}\colon \Lambda_Y
\to K_Y$, on obtient un diagramme commutatif dans $\detaleplus X$.
\[
\xymatrix{
\ar[d]_{f^\star(\Cl {Y/S})}f^\star \Lambda_Y\ar[r] & f^!\Lambda_Y
   \ar[d]^{f^!(\Cl {Y/S})} \\
f^\star K_Y\ar[r] & f^!K_Y }
\]

\emph{Via} l'isomorphisme canonique $f^\star\Lambda_Y\simeq \Lambda_X$, le
morphisme du haut s'identifie au morphisme
$\Cl f\colon\Lambda_X\to f^!\Lambda_Y$ ;
celui de droite est $f^!(\Cl {Y/S})$. D'après le
théorème~\ref{16-theoreme-cl-f-transitivite}, il vient que le morphisme
composé
$\Lambda_X\simeq f^\star \Lambda_Y\to f^!K_Y\simeq K_X$ est le morphisme de
Gysin $\Cl{X/S}$. On déduit de ceci un diagramme commutatif de la forme
suivante dans $\detaleplus X$ :

\[
\xymatrix{
\ar[d]_{f^\star(\Cl {Y/S})}f^\star \Lambda_Y\ar[r]^\sim &
\Lambda_X\ar[d]^{\Cl{X/S}} \ar[r] & f^!\Lambda_Y\ar[d]^{f^!\Cl{Y/S}} \\
f^\star K_Y\ar[r] & K_X\ar[r]^\sim & f^!K_Y}
\]

Comme $f$ est propre, on obtient par adjonction un nouveau diagramme
commutatif dans $\detaleplus Y$ :
\[
\xymatrix{
\Lambda_Y\ar[d]_{\Cl{Y/S}} \ar[r] & \R f_\star
\Lambda_X\ar[d]^{\R f_\star (\Cl{X/S})}\ar[r] &\Lambda_Y 
\ar[d]^{\Cl{Y/S}}\\
K_Y\ar[r] & \R f_\star K_X\ar[r]& K_Y}
\]

Le diagramme ci-dessus met en évidence une relation entre les morphismes
$\Cl{Y/S}$ et $\R f_\star(\Cl{X/S})$. Comme va le montrer le lemme
suivant, le premier morphisme est un facteur direct du second, ce qui
montre que la pureté ponctuelle de $X$ implique celle de $Y$, achevant la
démonstration de la
proposition~\ref{16-proposition-purete-argument-de-transfert}.

\begin{lemme2}
Sur le diagramme précédent,
les morphismes composés $\Lambda_Y\to \Lambda_Y$ et
$K_Y\to K_Y$ sont les multiplications par le degré $d$ (en particulier, ce
sont des isomorphismes).
\end{lemme2}

Comme $Y$ est connexe (non vide), on a un isomorphisme évident 
$\Lambda\iso \End_{\detaleplus Y}(\Lambda_Y)$. D'après le théorème de bidualité
locale (cf. \sgaquatreetdemi{Th.~finitude}{4.3}),
on a aussi un isomorphisme $\Lambda\iso
\End_{\detaleplus Y}(K_Y)$. Il suffit donc d'obtenir la conclusion au-dessus d'un
ouvert non vide de $Y$.
Quitte à remplacer $Y$ par un ouvert non vide convenable, on peut supposer
que $f$ est un revêtement étale.
On est ainsi ramené au lemme suivant :

\begin{lemme2}
Soit $f\colon X\to Y$ un morphisme de schémas fini étale de degré constant
$d$. Pour tout objet $M\in \detaleplus Y$, le morphisme composé
\[M\to f_\star f^\star M\to M\]
déduit des adjonctions canoniques $(f^\star,f_\star)$ et
$(f_\star,f^\star)$ est la multiplication par $d$.
\end{lemme2}

Grâce aux formules de projection, on peut supposer que $M=\Lambda_Y$. Il
suffit alors d'établir le résultat après un changement de base étale (non
vide) trivialisant le revêtement $X\to Y$ (par exemple une clôture
galoisienne de ce revêtement). Bref, on peut supposer que $X$ est une
réunion disjointe de $d$ copies de $Y$, auquel cas le résultat est trivial.

\bigskip

\begin{definition2}\label{16-definition-schemas-v}
Soit $(e_1,\dots,e_n)\in\NN^n$.
On définit un $S$-schéma :
\[V(S,\pi,e_1,\dots,e_n)=
\SP\left(\OO_S[T_1,\dots,T_n]/
\left(\textstyle\prod_{i=1}^n T_i^{e_i}-\pi\right)\right)\;\text{.}\]
Pour tout $i$, on note $H_i$ le sous-schéma fermé de
$V(S,\pi,e_1,\dots,e_n)$ défini par l'équation $T_i=0$.
\end{definition2}

\begin{proposition2}\label{16-proposition-standards-reguliers}
Soit $(e_1,\dots,e_n)$ un $n$-uplet d'entiers naturels non tous nuls.
Alors,
le schéma $V(S,\pi,e_1,\dots,e_n)$ est régulier et ponctuellement pur.
\end{proposition2}

On peut supposer que l'anneau des coefficients est $\ZZ/\ell\ZZ$ où $\ell$
est un nombre premier inversible sur $S$.

\begin{lemme2}~

\begin{enumerate}
\item[(i)] Si $(e_1,\dots,e_n)$ est un $n$-uplet d'entiers non tous nuls
dont au moins un est inversible dans $\eta$, le $S$-schéma
$V(S,\pi,e_1,\dots,e_n)$ est intègre, régulier et de fibre générique
lisse ;
\item[(ii)] Soit $d\geq 1$, si $S'$ est le trait obtenu en
extrayant une racine $d$-ième $\pi'$
de l'uniformisante $\pi$\;\footnote{Rappelons brièvement pourquoi $S'$ est
bien un trait. Si $A$ est un anneau de valuation discrète, $\pi$ une
uniformisante et $d\geq 1$, notons $A'=A[X]/(X^d-\pi)$. Il est évident que
$A'$ est local noethérien et que son idéal maximal est
engendré par $X$. On vérifie aussitôt que la
suite évidente $0\to A'\vers X A' \to A'/(X) \to 0$ de $A'$-modules est
exacte, ce qui donne une résolution projective du corps résiduel de $A'$.
D'après \cite[§D.1, chapitre~IV]{Algebre@Serre},
$A'$ est régulier de dimension
$0$ ou $1$ et comme $A'$ n'est pas un corps, c'est un anneau de valuation
discrète.},
pour tout $n$-uplet $(e_1,\dots,e_n)$, on a un
isomorphisme de schémas
\[V(S',\pi',e_1,\dots,e_n)=V(S,\pi,de_1,\dots,de_n)\;\text{;}\]
\item[(iii)] Si $(e_1,\dots,e_n)$ est un $n$-uplet d'entiers non tous nuls,
le schéma $V(S,\pi,e_1,\dots,e_n)$ est régulier et intègre ;
\item[(iv)] Soit $(e_1,\dots,e_n)$ un $n$-uplet d'entiers non tous nuls,
le schéma $V(S,\pi,e_1,\dots,e_n)$ est ponctuellement pur si et seulement si
pour tout $i$ tel que $e_i>0$, le couple régulier
$(V(S,\pi,e_1,\dots,e_n),H_i)$ est pur\;\footnote{Si
$e_i=0$, c'est vrai aussi : c'est un cas particulier du théorème de pureté
relatif, cf.~\sga{4}{xvi}{3.7}.} ;
\item[(v)] Soit $(e_1,\dots,e_n)$ un $n$-uplet d'entiers non nuls,
soit $e$ le p.p.c.m. des $e_i$ ;
on suppose que $\ell$ ne divise pas $e$ ; si
$V(S,\pi,e,\dots,e)$ est
ponctuellement pur, alors $V(S,\pi,e_1,\dots,e_n)$ aussi ;
\item[(vi)] Soit $(e_1,\dots,e_n)$ un $n$-uplet d'entiers non tous nuls,
$V(S,\pi,e_1,\dots,e_n)$ est ponctuellement pur si et seulement si
$V(S,\pi,\ell e_1,e_2,\dots,e_n)$ est ponctuellement pur.
\end{enumerate}
\end{lemme2}

Les assertions (i) et (ii) sont laissées en exercice au lecteur.
L'assertion (iii) résulte aussitôt de (i) et de (ii).

Pour montrer l'assertion~(iv),
il suffit d'observer que les diviseurs $H_i$ pour
$e_i>0$ sont
ponctuellement purs (ce sont des espaces affines sur le corps résiduel de
$S$) et forment un recouvrement de la fibre spéciale de
$V(S,\pi,e_1,\dots,e_n)$. La fibre générique du schéma
$V(S_,\pi,e_1,\dots,e_n)$
étant ponctuellement pure (puisque lisse sur une extension
de $\eta$), on peut conclure en
utilisant la proposition~\ref{16-proposition-deux-sur-trois}.

Concernant l'assertion~(v),
l'élévation des $T_i$ à la puissance $\frac e
{e_i}$ définit un morphisme fini et plat $V(S,\pi,e,\dots,e)\to
V(S,\pi,e_1,\dots,e_n)$ de degré $\frac{e^n}{e_1\dots e_n}$ (premier à
$\ell$) ; compte tenu du critère~(iv), la
proposition~\ref{16-proposition-purete-argument-de-transfert-facile} permet de
conclure.

Pour établir~(vi), remarquons que l'élévation de $T_1$ à la puissance
$\ell$ définit un morphisme fini et plat 
$V(S,\pi,\ell e_1,e_2,\dots,e_n)\to V(S,\pi,e_1,\dots,e_n)$
de degré $\ell$ et étale en dehors du lieu d'annulation de $T_1$.
Il suffit
donc de montrer que $(V(S,\pi,\ell e_1,e_2,\dots,e_n),H_1)$
est pur si et seulement si
$(V(S,\pi,e_1,\dots,e_n),H_1)$ l'est, ce qui résulte
de la proposition~\ref{16-proposition-up-and-down}.

\medskip

Établissons la proposition~\ref{16-proposition-standards-reguliers}. D'après
l'assertion (iii), les schémas considérés sont réguliers. Pour établir leur
pureté ponctuelle, d'après
le corollaire~\ref{16-corollaire-lisse-sur-ponctuellement-pur}, on peut
supposer qu'aucun des exposants $e_i$ n'est nul. Dans le cas où les tous
les entiers $e_i$ valent $1$, le résultat est établi dans 
\cite[theorem~1.4]{DWORK@Illusie} (voir aussi
\cite[Satz~2.21]{UEBER@Rapoport-Zink}). Grâce à l'utilisation d'un trait auxiliaire,
l'assertion~(ii) permet d'en déduire
que pour tout entier $d\geq 1$, $V(S,\pi,d,\dots,d)$
est ponctuellement pur. En utilisant l'assertion~(v), on obtient que
$V(S,\pi,e_1,\dots,e_n)$ est ponctuellement pur si $\ell$ ne divise aucun
des entiers $e_i$. L'assertion~(vi) permet de passer au cas général.

%>>>2
\subsection{Géométrie logarithmique}%<<<2
\label{16-subsection-geometrie-logarithmique}

\begin{definition2}
Soit $S$ un trait, de point générique $\eta$. La log-structure canonique
sur $S$ est la log-structure image directe de la log-structure triviale sur
$\eta$. Toute uniformisante de $S$ définit un morphisme de monoïdes
$\NN\to\Gamma(S,\OO_S)$ donnant naissance à une carte 
$S\to \SP(\ZZ[\NN])$ du log-schéma $S$.
\end{definition2}

L'objectif de cette sous-section d'établir le théorème suivant :

\begin{theoreme2}\label{16-theoreme-log-lisse-regulier-implique-p-pur}
Soit $S$ un trait muni de sa log-structure canonique. Soit $(X,M)\to S$  un
morphisme log-lisse de log-schémas fs. Si le schéma $X$ est régulier, alors
il est ponctuellement pur.
\end{theoreme2}

La proposition suivante précise \cite[theorem~3.5]{LSFI@Kato} dans le cas des
log-schémas fs :

\begin{proposition2}\label{16-proposition-critere-log-lissite-fs}
Soit $(X,M)\to (Y,N)$ un morphisme log-lisse entre log-schémas fs. On
suppose donnée une carte $Y\to\SP(\ZZ[Q])$ de $(Y,N)$ où $Q$ est un monoïde
fs sans torsion\;\footnote{Si $\overline y$ est un point géométrique de
$Y$, il existe un voisinage étale de $\overline y$ admettant une telle
carte avec $Q=M_{\overline y}/\OO_{Y,\overline y}^\times$ qui est fs 
saillant.}.
Pour tout point géométrique $\overline x$ de $X$, il
existe un voisinage étale $U$ de $\overline x$, un morphisme injectif de
monoïde $Q\to P$ avec $P$ fs sans torsion et une carte $U\to \ZZ[P]$ tels
que la partie de torsion de $\coker(Q^{\gp}\to P^{\gp})$ soit 
d'ordre inversible sur $U$ et que le morphisme de schémas
$U\to Y\times_{\SP(\ZZ[Q])}\SP(\ZZ[P])$ soit étale.
\end{proposition2}

Dans la démonstration du critère de log-lissité de
\cite[theorem~3.5]{LSFI@Kato}, des éléments $t_1,\dots,t_r$ de $M_{\overline x}$
sont choisis de sorte que la famille 
$(d\log t_1,\dots,d\log t_r)$ forme une base du faisceau des
log-différentielles
$\omega^1_{X/Y,\overline x}$. On considère ensuite le morphisme de monoïdes
évident $\NN^r\oplus Q\to M_{\overline x}$ donné sur la composante $\NN^r$
par les $t_1,\dots,t_r$. Il est tel que le conoyau de $\ZZ^r\oplus
Q^{\gp}\to M_{\overline x}^{\gp}$ soit fini d'exposant $n$ inversible dans
l'anneau $\OO_{X,\overline x}$ (en particulier, $\OO_{X,\overline
x}^\times$ est $n$-divisible). Il existe un
morphisme injectif $\ZZ^r\oplus Q^{\gp}\to G$ de conoyau tué par une
puissance de $n$ et un
prolongement $h\colon G\to M_{\overline x}^{\gp}$
de $\ZZ^r\oplus Q^{\gp}\to M_{\overline x}^{\gp}$
tel que $G\to M_{\overline
x}^{\gp}/\OO^\times_{X,\overline x}$ soit surjectif. Comme $M_{\overline
x}^{\gp}/\OO^\times_{X,\overline x}$ est un groupe abélien de type fini et
sans torsion, le lemme suivant montre que l'on peut s'arranger pour que $G$
soit un groupe abélien libre.
Dans la démonstration de \cite[theorem~3.5]{LSFI@Kato}, on pose ensuite
$P=h^{-1}(M_{\overline x})$ et il est montré que sur un voisinage étale $U$
de $\overline x$, $P$ engendre la log-structure de $(X,M)$ et que
le morphisme
de schémas
$U\to S\times_{\SP(\ZZ[Q])}\SP(\ZZ[P])$ est étale en $\overline x$. Le
monoïde $P$ ainsi construit est fs et sans torsion.

\begin{lemme2}
Soit $n$ un entier naturel non nul.
Soit $A$ un groupe abélien libre de type fini.
Soit $\varphi\colon A\to B$ un morphisme de groupes abéliens.
Soit $U\subset B$ un sous-groupe $n$-divisible. On suppose que $B/U$ est
sans torsion et que $\coker(A\to
B/U)$ est fini et tué par $n$.
Alors, il existe un groupe abélien $A'$ libre de type fini,
un morphisme injectif $A\to A'$ de groupes abéliens tel que $A'/A$ soit tué
par une puissance de 
$n$ et une extension $A'\to B$ du morphisme $A\to B$ telle que le
morphisme composé $A'\to B/U$ soit surjectif.
\end{lemme2}

Grâce à une récurrence sur l'ordre de $\coker(A\to B/U)$, on peut supposer
que $\coker(A\to B/U)$ est cyclique d'ordre $d\geq 2$, engendré par la
classe d'un élément $b\in B$. Il existe donc $a\in A$ et $u\in U$ tels que
$db=\varphi(a)+u$. Comme $u$ est $n$-divisible, il existe $\tilde{u}\in U$
tel que $u=d\tilde{u}$. Quitte à remplacer $b$ par $b-\tilde{u}$, on peut
supposer que $u=0$. On forme le carré cocartésien suivant dans la catégorie
des groupes abéliens :

\[
\xymatrix{ \mathbf{Z}\ar[r]^a\ar[d] & A \ar[d]\\
\frac 1 d \mathbf{Z}\ar[r] & A'}
\]

En raison de la relation $db=\varphi(a)$, on peut définir un unique
morphisme de groupes abéliens $\varphi'\colon A'\to B$ induisant
$\varphi\colon A\to B$ et envoyant $\frac 1 d$ sur $b$. On obtient ainsi
une surjection $A'\to B/U$ induisant un isomorphisme $A'/A\iso
\coker(A\to B/U)$. Il reste à vérifier que $A'$ est sans torsion. Soit $a'$
un élément de torsion de $A'$. L'image de $a'$ dans $B/U$ \emph{via}
$\varphi'$ est de torsion, mais $B/U$ étant sans torsion, on a
$\varphi(a')\in U$. Comme $\varphi'$ induit un isomorphisme $A'/A\iso
\coker(A\to B/U)$, on en déduit que $a'\in A$, mais $A$ est sans torsion,
donc $a'=0$.

\begin{proposition2}\label{16-proposition-structure-locale-log-lisse-regulier}
Soit $(X,M)\to S$ un log-schéma fs log-lisse sur un trait $S$ (muni de sa
log-structure canonique). On suppose que le schéma $X$ est régulier. Alors,
localement pour la topologie étale, $X$ admet un morphisme étale vers un
schéma $V(S,\pi,e_1,\dots,e_n)$ où $(e_1,\dots,e_n)$ est un $n$-uplet
d'entiers non tous nuls (cf.~définition~\ref{16-definition-schemas-v}).
\end{proposition2}

Soit $\pi$ une uniformisante de $S$ ; elle donne naissance
à une carte $S\to \SP(\ZZ[\NN])$.
D'après la
proposition~\ref{16-proposition-critere-log-lissite-fs}, on peut supposer qu'il
existe un monoïde $P$ fs sans torsion, un morphisme injectif $\NN\to P$, une
carte $X\to\SP(\ZZ[P])$ telle que le morphisme de schémas 
$X\to S\times_{\SP(\ZZ[\NN])}\SP(\ZZ[P])$ soit étale. Soit $\overline x$ un
point géométrique de $X$. On note $P'$ le sous-monoïde de $P$ formé des
éléments dont l'image dans $\Gamma(X,\OO_X)$ soit inversible au point
$\overline x$.

On peut supposer que $P'$ est un groupe. En effet, si $A$
est un sous-ensemble fini de $P'$ qui engendre le
groupe abélien (libre de type fini) $P'^{\gp}$\;\footnote{En fait,
on peut montrer que $P'$ est un monoïde de type fini
(c'est une face de $P$).},
on peut remplacer $X$ par
le voisinage ouvert de $\overline x$ sur lequel les images des éléments de
$A$ (et donc de $P'$) sont inversibles dans le faisceau structural et par
suite, remplacer $P$ par $P[-P']$ qui est encore fs et sans torsion.

Le fait que $X\to \SP(\ZZ[P])$ soit une carte implique alors que $P'$
est le noyau de
$P^{\gp}\to M_{\overline x}^{\gp}/\OO^\times_{X,\overline x}$. En
particulier, on obtient un isomorphisme
\[P/P'\iso M_{\overline
x}/\OO^\times_{X,\overline x}\;\text{.}\]
Comme $X$ est log-régulier, on reconnaît que $X$ est régulier au fait que
$M_{\overline x}/\OO^\times_{X,\overline x}$ soit un monoïde libre. Par
conséquent, il existe un entier $r$ et un isomorphisme de monoïdes
$\NN^r\iso P/P'$. On peut relever ce morphisme en un morphisme $\NN^r\to
P$, ce qui permet de construire un isomorphisme $\NN^r\oplus P'\iso
P$.

Il en résulte que le morphisme de carte $X\to \SP(\ZZ[P])$ a pour but un
schéma isomorphe à
$\SP(\ZZ[\NN^r\oplus P'])$ qui est le produit d'un espace affine
et d'un tore déployé (dont $P'$ est le groupe des caractères). Dans la
carte du morphisme $(X,M)\to S$ qui est donnée, l'image de $1$ par le
morphisme de monoïdes $\NN\to P$ peut s'écrire $(e_1,\dots,e_r,p')$ dans
$\NN^r\oplus P'$ \emph{via} les identifications ci-dessus. On peut choisir
une base $a_1,\dots,a_s$ de $P'$ comme groupe abélien telle que
$p'=\sum_{i=1}^s f_i a_i$ avec $f_i\in\NN$. On a ainsi construit un
morphisme étale $X\to V(S,\pi,e_1,\dots,e_r,f_1,\dots,f_s)$ (avec les
$e_1,\dots,e_r,f_1,\dots,f_s$ non tous nuls).

\medskip

Compte tenu de la proposition~\ref{16-proposition-standards-reguliers},
le théorème~\ref{16-theoreme-log-lisse-regulier-implique-p-pur} résulte
aussitôt de la
proposition~\ref{16-proposition-structure-locale-log-lisse-regulier}.


%>>>2
\subsection{Démonstration du théorème de pureté}%<<<2
\label{16-subsection-demonstration-theoreme-purete}

Démontrons le théorème~\ref{16-theoreme-purete-absolue}.
D'après les propositions~\ref{16-proposition-lemme-chinois} et
\ref{16-proposition-ops-n-egale-ell}, on peut supposer que l'anneau des
coefficients $\Lambda$ est $\ZZ/\ell\ZZ$ où $\ell$ est un nombre premier.
D'après la proposition~\ref{16-proposition-deux-sur-trois}, il s'agit de montrer
que tout $\ZZn[\ell]$-schéma régulier est ponctuellement pur. D'après
\cite[corollary~6.1.5]{AZUMINO@Fujiwara}, on peut supposer que $X$ est un schéma
régulier intègre, quasi-projectif et plat 
sur un trait (strictement hensélien) $S$, que
l'on peut supposer d'inégale caractéristique d'après
la proposition~\ref{16-proposition-purete-egale-caracteristique}. On peut
utiliser les notations de la
sous-section~\ref{16-subsection-schemas-sur-un-trait}.
Quitte à étendre le
trait $S$, on peut supposer que l'anneau sous-jacent à $S$ est
intégralement fermé dans le corps des fonctions rationnelles sur $X$. La
fibre générique $X_\eta$ de $X$ est donc géométriquement intègre.

En appliquant \cite[proposition~4.4.1]{BS@Vidal} à la normalisation de
l'adhérence de $X$ dans un plongement projectif, on obtient qu'il existe un
groupe fini $G$ et un diagramme $G$-équivariant :
\[
\xymatrix{X'\ar[r]\ar[d] & \ar[d] X \\ S'\ar[r] & S}
\]
tels que :
\begin{itemize}
\item $G$ agisse trivialement sur $X$ et $S$ ;
\item $S'\to S$ soit une extension finie de traits ;
\item $X'\to X$ soit projectif, $X'$ soit régulier, connexe
et à réduction semi-stable sur $S'$ ;
\item $G$ agisse fidèlement sur $X'$ et $X'\to X$ soit génériquement un
revêtement étale galoisien de groupe $G$.
\end{itemize}

On munit $X'$ de la log-structure dont l'ouvert de trivialité est la fibre
générique de $X'\to S'$. Soit $H$ un $\ell$-Sylow de $G$. On note $T=S'/H$.
L'extension de traits (strictement henséliens) $S'\to T$ est d'ordre une
puissance de $\ell$, donc modérément ramifiée. Par conséquent, pour les
log-structures canoniques, $S'\to T$ est log-étale. Comme on sait que $X'$
est log-lisse sur $S'$, il l'est donc aussi sur $T$. Comme $H$ agit
trivialement sur $T$ et que son action sur $X'$ est modérée, on peut
appliquer le théorème de résolution équivariante qui
donne un morphisme projectif et birationnel $H$-équivariant
$X''\to X'$ de log-schémas tel que $X''$ soit log-lisse sur $T$ et que $H$
agisse très modérément sur $X''$. Le log-schéma quotient $X''/H$ est donc
log-lisse sur $T$ (en particulier, $X''/H$ est 
log-régulier). D'après le théorème de résolution
des log-singularités, il existe un log-éclatement
(en particulier, log-étale,
projectif et birationnel) $X'''\to X''/H$ tel que $X'''$
soit régulier. La situation est résumée sur le diagramme suivant :

\[
\xymatrix{& \ar@(lu,ru)^{H} X''\ar[r]^{\text{birat.}}\ar[d] 
& \ar@(lu,ru)^{G}X'\ar[r]\ar[d] & X\ar[dd] \\
X'''\ar[r]^{\text{birat.}}&\ar@{.>}[rru] X''/H\ar[rd] & \ar@(dl,ul)^{G}S'\ar[d] & \\
& & T\ar[r] &S}
\]

Le log-schéma $X'''$ est régulier et log-lisse sur $T$ ;
d'après le théorème~\ref{16-theoreme-log-lisse-regulier-implique-p-pur},
$X'''$ est ponctuellement pur. Le morphisme évident $X'''\to X$ est
projectif et génériquement un revêtement
étale de degré premier à $\ell$, d'après la 
proposition~\ref{16-proposition-purete-argument-de-transfert}, on peut
conclure que $X$ est ponctuellement pur, ce qui achève la démonstration du
théorème de pureté.



%>>>2

\ifx\danslelivre\undefined
\bibliography{../fichiers-configuration-latex/bibliographie}
\bibliographystyle{../fichiers-configuration-latex/smfalpha}
\end{document}
\fi
